\begin{exname}Problem\end{exname}
% begin pdf problem style
\renewcommand{\ocExName}{Problem}
% end pdf problem style
\chapter{Sums of random pieces}\label{ch:harder}

\begin{center}\itshape
Add up many small random pieces and a bell curve appears.\\
How soon? Point by point, or only on intervals? What if the pieces depend on each other,\\
or if there are infinitely many of them? Each question has an answer, and each answer has a price.
\end{center}

Harald Cram\'er met Jarl Waldemar Lindeberg in 1922, the year Lindeberg published the condition that carries his name and the short proof of the
central limit theorem that every probabilist still learns. Lindeberg owned a beautiful farm in Finland. In his recollections of 1976 Cram\'er tells
that when colleagues reproached Lindeberg for not being active enough in research, he answered ``Well, I am really a farmer'', and that when someone
remarked that his farm was not properly cultivated, he answered that his real job was to be a professor.

The proof he found that year is a farmer's proof: it replaces the random pieces one at a time by Gaussian pieces of the same mean and variance, and
counts the damage of each swap with a Taylor expansion. Nothing in it is specific to coin tosses or to sums of identical terms, which is why it keeps
coming back. Volume~11 sharpened the theorem in every direction: Lindeberg's condition is necessary (Feller), the error is of order $n^{-1/2}$ (Berry and
Esseen), the probability of each single value is Gaussian too (the local limit theorem), Stein's identity measures the distance even for dependent
pieces, and Kolmogorov's three series decide when an infinite sum converges at all.

This chapter takes those tools to places where nobody says they apply. The first section finds independence hidden in permutations, tilings, set
partitions and the binary digits of a random number. The second asks how fast the limit is reached, for fixed points, prime factors, arithmetic
progressions, ranks and planar walks. The third counts exactly, with the local limit theorem: lucky tickets, cancelling signs, two samples
compared by Smirnov. The last adds infinitely many pieces and asks what the sum looks like: smooth, singular or a fractal dust, with a famous answer of
Erd\H os for the golden ratio and Kac's count of the real roots of a random polynomial.

\begin{think}
Choose signs at random in $\pm1\pm2\pm3\cdots\pm n$. The sum is $0$ only if $\frac{n(n+1)}2$ is even. When it is, roughly how many of the $2^n$ choices
give $0$: about $\frac{2^n}n$, about $\frac{2^n}{n^{3/2}}$, or about $\frac{2^n}{n^2}$? Guess before Problem~\ref{pb:signedsums}.
\end{think}

Throughout, $Z$ is a standard Gaussian variable, $\phi(x)=\frac1{\sqrt{2\pi}}e^{-x^2/2}$ and $\Phi$ its density and distribution function, and
$\varepsilon_1,\varepsilon_2,\dots$ are independent fair signs ($\pm1$ with probability $\frac12$ each). For two laws, $d_K$ is the Kolmogorov distance
(the largest gap between distribution functions), $d_W$ the Wasserstein distance and $\|\cdot\|$ the total variation distance, as in Volume~11, Chapter~6.

\section{Independence in disguise}

Lindeberg's theorem wants independent pieces. Very often the pieces are not independent as given, but become independent, or nearly so, after a change of
point of view: a uniform number in $[0,1]$ is a sequence of independent digits, a random permutation is a sequence of independent uniform numbers, a
real-rooted polynomial is a product of independent coins.

\begin{exercise}[$\star$]\label{pb:pairwisebits}
\emph{Pairwise is not enough.} For a finite non-empty set $S$ of positive integers let $\chi_S=\prod_{i\in S}\varepsilon_i$.
\begin{enumerate}[label=(\alph*)]
\item Show that the $\chi_S$ are fair signs and pairwise independent, and that $\chi_{\{1,2\}}$, $\chi_{\{1\}}$, $\chi_{\{2\}}$ are not independent.
\item List the non-empty finite sets $S$ in the order of $\sum_{i\in S}2^{i-1}$, as $Y_1,Y_2,Y_3,\dots$, and let $T_N=Y_1+\dots+Y_N$. Show that $T_N\in\{-1,2^k-1\}$ for
$N=2^k-1$, and that $\frac{T_N}N\to0$ almost surely along all $N$.
\item Show that $\frac{T_N}{\sqrt N}$ does not converge in law to $N(0,1)$ along $N=2^k-1$, although each $Y_n$ has mean $0$ and variance $1$.
\item A Monte Carlo method estimates $\mathbb Ef(U)$ by averaging $f$ at $N$ random points. Explain why $k$ random bits can replace $N=2^k-1$ independent
samples in Chebyshev's error bound, but not in a Gaussian confidence interval.
\end{enumerate}
\end{exercise}

\begin{exercise}[$\star\star$]\label{pb:sawtooth}
\emph{The variance that is not $\frac1{12}$.} Let $B(x)=\{x\}-\frac12$ be the sawtooth, and $x$ uniform in $[0,1]$.
\begin{enumerate}[label=(\alph*)]
\item Show that $\int_0^1B(x)B(mx)\dd x=\frac1{12m}$ for every integer $m\ge1$.
\item Show that $\frac1{\sqrt n}\sum_{k=0}^{n-1}B(2^kx)$ converges in law to $N(0,\sigma^2)$, and compute $\sigma^2$. Compare with $\int_0^1B^2=\frac1{12}$, the variance one
would guess if the terms were independent.
\item Do the same for $\frac1{\sqrt n}\sum_{k<n}B(3^kx)$.
\item Show that $\sigma^2=\int_0^1B^2+2\sum_{k\ge1}\int_0^1B(x)B(2^kx)\dd x$ in (b), and find a non-zero trigonometric polynomial $f$ of mean $0$ for which
$\frac1{\sqrt n}\sum_{k<n}f(2^kx)\to0$ in probability.
\end{enumerate}
\end{exercise}

\begin{exercise}[$\star\star$]\label{pb:peaks}
\emph{Peaks.} A \emph{peak} of a permutation $\pi$ of $\{1,\dots,n\}$ is a position $1<i<n$ with $\pi(i-1)<\pi(i)>\pi(i+1)$. Let $P_n$ be the number of peaks of a
uniformly random permutation.
\begin{enumerate}[label=(\alph*)]
\item Show that $\mathbb EP_n=\frac{n-2}3$ and $\operatorname{Var}P_n=\frac{2(n+1)}{45}$ for $n\ge4$.
\item Let $X_1,X_2,\dots$ be centred, bounded by $c$, and \emph{$m$-dependent}: $(X_i)_{i\le j}$ is independent of $(X_i)_{i>j+m}$ for every $j$. Let
$S_n=X_1+\dots+X_n$, and suppose $\operatorname{Var}S_n\ge\delta n$ for some $\delta>0$ and all large $n$. Show that $\frac{S_n}{\sqrt{\operatorname{Var}S_n}}\Rightarrow N(0,1)$.
\item Deduce that $\bigl(P_n-\frac{n-2}3\bigr)\big/\sqrt{\frac{2n}{45}}\Rightarrow N(0,1)$.
\end{enumerate}
\end{exercise}

\begin{exercise}[$\star\star$]\label{pb:tilings}
\emph{Random tilings.} A strip of length $n$ is tiled with squares of length $1$ and dominoes of length $2$; there are $F_{n+1}$ tilings, where $F$ is the
Fibonacci sequence. Choose one uniformly, and let $D_n$ be its number of dominoes.
\begin{enumerate}[label=(\alph*)]
\item Let $T_1,T_2,\dots$ be independent tiles, a domino with probability $q=\varphi^{-2}$ and a square with probability $1-q=\varphi^{-1}$, where
$\varphi=\frac{1+\sqrt5}2$, laid end to end from $0$. Show that, conditioned on the event $H_n$ that some tile ends exactly at $n$, the tiling of $[0,n]$ is uniform.
\item Show that $P(H_n)=\frac{1-(-q)^{n+1}}{1+q}$.
\item Show that $\frac{D_n-an}{\sqrt n}\Rightarrow N(0,s^2)$, with $a=\frac{5-\sqrt5}{10}$ and $s^2=\frac1{5\sqrt5}$.
\end{enumerate}
\end{exercise}

\begin{exercise}[$\star\star\star$]\label{pb:harper}
\emph{Harper's trick.} Let $K_n$ be the number of blocks of a uniformly random partition of $\{1,\dots,n\}$, so that $P(K_n=k)=\frac{S(n,k)}{B_n}$, with the
Stirling numbers of the second kind $S(n,k)$ and the Bell number $B_n=\sum_kS(n,k)$.
\begin{enumerate}[label=(\alph*)]
\item Let $X$ take values in $\{0,\dots,n\}$, and suppose that $G(x)=\sum_kP(X=k)x^k$ has only real roots. Show that $X$ has the law of a sum of
$n$ independent Bernoulli variables.
\item Show that $A_n(x)=\sum_kS(n,k)x^k$ satisfies $A_{n+1}(x)=x\bigl(A_n(x)+A_n'(x)\bigr)$, and deduce that $A_n$ has $n$ distinct real roots, all in $]-\infty,0]$.
\item Let $M$ be an integer variable with $P(M=m)=\frac{m^n}{e\,m!\,B_n}$ (Dobi\'nski's formula, Volume~5, Exercise~2.3, says this is a law). Show that
$\mathbb EK_n=\mathbb EM-1$ and $\operatorname{Var}K_n=\operatorname{Var}M-1$.
\item Show that $\operatorname{Var}M\to\infty$, and deduce that $\frac{K_n-\mathbb EK_n}{\sqrt{\operatorname{Var}K_n}}\Rightarrow N(0,1)$.
\end{enumerate}
\end{exercise}

\section{How fast?}

A limit theorem says that an error tends to $0$; a rate says how large it is now. The rates of Volume~11 come in three flavours: Berry--Esseen for
Kolmogorov's distance, Stein's method for Wasserstein's distance and for dependent pieces, Chen--Stein for Poisson limits. Each problem below needs one
of them, and sometimes an exact formula does better than all three.

\begin{exercise}[$\star$]\label{pb:fixedtv}
\emph{Fixed points, superexponentially.} Let $F_n$ be the number of fixed points of a uniformly random permutation of $\{1,\dots,n\}$.
\begin{enumerate}[label=(\alph*)]
\item Show that $\bigl\|\mathcal L(F_n)-\text{Poisson}(1)\bigr\|\le\frac{2^n}{(n+1)!}$.
\item Compare with what the Chen--Stein method gives when $F_n$ is written as a sum of $n$ indicators.
\end{enumerate}
\end{exercise}

\begin{exercise}[$\star$]\label{pb:kubslow}
\emph{A very slow Gaussian.} Let $B_p$, for primes $p\le x$, be independent Bernoulli variables with $P(B_p=1)=\frac1p$, and $W_x=\frac{\sum_p(B_p-1/p)}{\sigma_x}$ with
$\sigma_x^2=\sum_{p\le x}\frac1p(1-\frac1p)$. This is the model of the number of prime factors of a random integer used in Volume~6, Chapter~9.
\begin{enumerate}[label=(\alph*)]
\item Show that $d_W(W_x,Z)\le\frac3{\sigma_x}$, and deduce that $d_K(W_x,Z)\le\frac{2}{(\log\log x)^{1/4}}$ for $x$ large.
\item How large must $x$ be for this bound on $d_K$ to fall below $1\%$?
\item Show that in fact $d_K(W_x,Z)\ge\frac c{\sqrt{\log\log x}}$ for some $c>0$ and all large $x$. (For $x=10^6$ one finds numerically $d_K(W_x,Z)=0.15$, while the bound of (a) is useless there.)
\end{enumerate}
\end{exercise}

\begin{exercise}[$\star\star$]\label{pb:twelve}
\emph{Twelve uniforms.} A cheap way to produce roughly Gaussian numbers is $W=U_1+\dots+U_{12}-6$, with $U_i$ independent and uniform on $[0,1]$.
\begin{enumerate}[label=(\alph*)]
\item Let $X_1,\dots,X_n$ be independent, symmetric, with $\mathbb EX_i^2=1$ and $\mathbb EX_i^4=m_4$, and $W=\frac{X_1+\dots+X_n}{\sqrt n}$. Show that for every $h$ with four
bounded continuous derivatives,
\[
\bigl|\mathbb Eh(W)-\mathbb Eh(Z)\bigr|\le\frac{(m_4+3)\,\|h^{(4)}\|_\infty}{24\,n}.
\]
\item Apply it to $W=U_1+\dots+U_{12}-6$: show that $|\mathbb E\cos W-e^{-1/2}|\le\frac1{60}$, and compute $\mathbb E\cos W$ exactly.
\item Show that $P(|W|>6)=0$ and $P(|W|>4)=\frac{2(2^{12}-12)}{12!}\approx1.7\times10^{-5}$, while $P(|Z|>4)\approx6.3\times10^{-5}$. What does (a) say, and not say, about tails?
\end{enumerate}
\end{exercise}

\begin{exercise}[$\star\star$]\label{pb:threeap}
\emph{Random sets without progressions.} Let $N$ be an odd prime and $A$ a random subset of $\Z/N\Z$ containing each residue independently with probability
$p=cN^{-2/3}$. A \emph{progression} is a set $\{a,a+d,a+2d\}$ with $d\ne0$.
\begin{enumerate}[label=(\alph*)]
\item Show that there are $\frac{N(N-1)}2$ progressions, and that two different progressions share at most two elements.
\item Show that the number $W$ of progressions contained in $A$ satisfies $\|\mathcal L(W)-\text{Poisson}(\lambda)\|=O(N^{-1/3})$, with $\lambda=p^3\frac{N(N-1)}2$.
\item Deduce that $P(A\text{ contains no progression})\to e^{-c^3/2}$, and that $|A|\sim cN^{1/3}$ with high probability.
\end{enumerate}
\end{exercise}

\begin{exercise}[$\star\star\star$]\label{pb:spearman}
\emph{Spearman's rank correlation, with a rate.} Let $b_1,\dots,b_n$ and $c_1,\dots,c_n$ be real with $\sum_ib_i=\sum_jc_j=0$ and $\bigl(\sum_ib_i^2\bigr)\bigl(\sum_jc_j^2\bigr)=n-1$,
let $\pi$ be a uniformly random permutation and $W=\sum_ib_ic_{\pi(i)}$.
\begin{enumerate}[label=(\alph*)]
\item Show that $\mathbb EW=0$ and $\operatorname{Var}W=1$. More generally, compute $\operatorname{Var}\sum_ix_iy_{\pi(i)}$ for arbitrary reals $x_i,y_j$.
\item Let $I\ne J$ be a uniformly random ordered pair of positions, independent of $\pi$, and $\pi'=\pi\circ(I\,J)$. Show that $(W,W')$, with
$W'=\sum_ib_ic_{\pi'(i)}$, is exchangeable and that $\mathbb E[W'-W\mid\pi]=-\frac2{n-1}W$.
\item Show that
\[
d_W(W,Z)\le\sqrt{\frac{\sum_ib_i^4\sum_jc_j^4}{n-1}}+\frac4n+\frac{16}{n-1}\sum_i|b_i|^3\sum_j|c_j|^3.
\]
\item \emph{Spearman's $\rho$} of a permutation is the correlation between $i$ and $\pi(i)$. Show that $\sqrt{n-1}\,\rho$ is a statistic $W$ as above, and that
$d_W(\sqrt{n-1}\,\rho,Z)=O(n^{-1/2})$.
\end{enumerate}
\end{exercise}

\begin{exercise}[$\star\star\star$]\label{pb:erdostaylor}
\emph{The returns of a planar walk.} Let $S$ be the simple random walk on $\Z^2$ started at $0$, $N_n$ the number of $t\in\{1,\dots,n\}$ with $S_{2t}=0$, and
$G_n=\mathbb EN_n$.
\begin{enumerate}[label=(\alph*)]
\item Show that $P(S_{2t}=0)=\bigl(\binom{2t}t4^{-t}\bigr)^2$, and that $G_n\sim\frac1\pi\log n$.
\item Show that $k!\,G_{\lfloor n/k\rfloor}^k\le\mathbb E[N_n(N_n-1)\cdots(N_n-k+1)]\le k!\,G_n^k$ for every $k\ge1$.
\item Deduce that $\frac{N_n}{G_n}$ converges in law to the exponential law of mean $1$. In particular $\frac{N_n}{\mathbb EN_n}$ does not tend to $1$: the walk comes home
about $\frac1\pi\log n$ times on average, but the actual number is as random as a waiting time.
\end{enumerate}
\end{exercise}

\section{Counting with the local limit theorem}

The local limit theorem (Theorem~4.12) turns a probability into a count: the number of objects with a given sum is the total number times
$\frac{1}{\sigma\sqrt{2\pi n}}$, up to $1+o(1)$. Its proof is a Fourier integral, and that integral is often an exact formula in disguise.

\begin{exercise}[$\star$]\label{pb:dice}
\emph{Dice modulo seven.} Let $S_n$ be the total of $n$ fair dice.
\begin{enumerate}[label=(\alph*)]
\item Show that $P(7\mid S_n)=\frac17+\frac67\bigl(-\frac16\bigr)^n$, and find $P(S_n\equiv r\pmod7)$ for every $r$.
\item Find $P(S_n\equiv r\pmod5)$ exactly. Show that modulo $4$ the distance to the uniform law decays like $\bigl(\frac{\sqrt2}6\bigr)^n$, and modulo $2$ it is $0$ for
every $n\ge1$.
\item A die has faces $1,4,7,10,13,16$. Show that its total $S_n$ satisfies $\sup_{k\equiv n\,(3)}\bigl|\sigma\sqrt n\,P(S_n=k)-3\phi\bigl(\frac{k-n\mu}{\sigma\sqrt n}\bigr)\bigr|\to0$,
with $\mu,\sigma$ to be found, and explain the factor $3$.
\end{enumerate}
\end{exercise}

\begin{exercise}[$\star$]\label{pb:squaresmodp}
\emph{Two squares make a uniform residue.} Let $p$ be an odd prime and $a_1,a_2,\dots$ independent and uniform in $\Z/p\Z$. Let $Q_k$ be the law of
$a_1^2+\dots+a_k^2$ modulo $p$ and $U$ the uniform law.
\begin{enumerate}[label=(\alph*)]
\item Show that $\|Q_k-U\|\le\frac12\sqrt{p-1}\,p^{-k/2}$ for every $k\ge1$.
\item Compute $\|Q_2-U\|$ exactly, and compare.
\item For the walk with steps $\pm1$ instead of squares, show that $\|Q_k-U\|$ stays above $\frac1{20}$ as long as $k\le p^2/10$, for $p$ large. Why are squares so
much better mixers?
\end{enumerate}
\end{exercise}

\begin{exercise}[$\star\star$]\label{pb:lucky}
\emph{Lucky tickets.} Call a ticket with a six-digit number \emph{lucky} when the sum of the first three digits
equals the sum of the last three. More generally, let $L_m$ be the number of strings of $2m$ decimal digits whose first half and second half have the same digit sum.
\begin{enumerate}[label=(\alph*)]
\item Show that $L_m=\frac1{2\pi}\int_{-\pi}^\pi\Bigl(\frac{\sin5t}{\sin\frac t2}\Bigr)^{2m}\dd t$.
\item Show that $L_m\sim\frac{10^{2m}}{\sqrt{33\pi m}}$.
\item There are $55\,252$ lucky six-digit tickets. Show that the next term of the expansion is
\[
L_m=\frac{10^{2m}}{\sqrt{33\pi m}}\Bigl(1-\frac{101}{1320\,m}+O\bigl(m^{-2}\bigr)\Bigr),
\]
and check that it brings the estimate for $m=3$ to within $0.02\%$ of the truth.
\end{enumerate}
\end{exercise}

\begin{exercise}[$\star\star\star$]\label{pb:signedsums}
\emph{Signs that cancel.} Let $Z(n)$ be the number of choices of signs with $\pm1\pm2\pm\dots\pm n=0$.
\begin{enumerate}[label=(\alph*)]
\item Show that $Z(n)=0$ unless $n\equiv0$ or $3\pmod4$, and that $Z(n)=\frac{2^n}{2\pi}\int_{-\pi}^\pi\prod_{k=1}^n\cos(kt)\dd t$.
\item Show that $Z(n)\sim\sqrt{\frac6\pi}\,\frac{2^n}{n^{3/2}}$ for $n\equiv0,3\pmod4$.
\item Deduce the number of subsets of $\{1,\dots,n\}$ whose sum is $\frac{n(n+1)}4$, and of ways to split $\{1,\dots,n\}$ into two parts of equal sum.
\end{enumerate}
\end{exercise}

\begin{exercise}[Journey]\label{pb:smirnov}
\emph{Smirnov's two samples.} Two samples of $n$ real numbers each, drawn independently from the same continuous law, are merged and sorted. Read the merged
list from left to right and step up for a number of the first sample, down for one of the second: this gives a path $S_0=0,S_1,\dots,S_{2n}=0$, uniformly
distributed among the $\binom{2n}n$ such paths. Smirnov's statistic is $D_n=\frac1n\max_k|S_k|$, the largest gap between the two empirical distribution functions.
\begin{enumerate}[label=(\alph*)]
\item Let $h\ge1$. Show that the number of paths with $2n$ steps $\pm1$, from $0$ to $0$, that never reach $h$ or $-h$, is
\[
\sum_{k\in\Z}(-1)^k\binom{2n}{n+kh}.
\]
\item Show that for each $k$, $\binom{2n}{n+k}\big/\binom{2n}n=e^{-k^2/n}\bigl(1+O(\frac{k^4}{n^3}+\frac{k^2}{n^2})\bigr)$ when $|k|\le n^{2/3}$, and bound it by $e^{-k^2/(2n)}$ for all $|k|\le n$.
\item Deduce that for every $x>0$,
\[
P\Bigl(\sqrt{\tfrac n2}\,D_n<x\Bigr)\longrightarrow K(x)=1-2\sum_{k\ge1}(-1)^{k-1}e^{-2k^2x^2}.
\]
\item Show that $K(x)=\frac{\sqrt{2\pi}}x\sum_{k\ge1}e^{-(2k-1)^2\pi^2/(8x^2)}$, and that one term of the left series, or two terms of the right one, give $K(1.36)=0.951$ to three digits. Which series would you use for $x=0.3$?
\item The same limit governs one sample compared with the uniform law (Kolmogorov, 1933; see Volume~11, Remark in Chapter~3). Direct computation gives
$\sqrt N\sup_t|F_N(t)-t|$ for the $N$ points $\{k\sqrt2\}$, $k\le N$: it is $0.03$ for $N=10^4$, far below the $0.5$ to $1.4$ that random points give. Would a goodness-of-fit
test accept the sequence $(\{k\sqrt2\})$ as random?
\end{enumerate}
\end{exercise}

\section{Infinitely many pieces}

Now the sum never stops. Kolmogorov's three series say whether it converges; once it does, its law can be smooth, or a dust, or something in between,
and the characteristic function, a product of infinitely many cosines, tells which. Two questions run through the section: is the law continuous, and is it
absolutely continuous?

\begin{exercise}[$\star$]\label{pb:primeseries}
\emph{Random signs on the primes.} Let $(\varepsilon_p)$ be independent fair signs indexed by the primes, and $a\in\R$.
\begin{enumerate}[label=(\alph*)]
\item Show that $\sum_p\frac{\varepsilon_p}{\sqrt p\,(\log p)^a}$ converges almost surely if $a>0$, and diverges almost surely if $a\le0$.
\item Show that $\sum_{n\ge2}\frac{\varepsilon_n}{\sqrt n\,(\log n)^a}$, over all integers, converges almost surely if and only if $a>\frac12$.
\item When $a\le0$ in (a), show that the partial sums satisfy $\limsup=+\infty$ and $\liminf=-\infty$ almost surely.
\item Deduce that the random Dirichlet series $\sum_p\varepsilon_pp^{-\sigma}$ converges almost surely for every $\sigma>\frac12$, and diverges almost surely for $\sigma=\frac12$.
\end{enumerate}
\end{exercise}

\begin{exercise}[$\star$]\label{pb:factorialseries}
\emph{A dust of dimension zero.} Let $X=\sum_{n\ge1}\frac{\varepsilon_n}{n!}$.
\begin{enumerate}[label=(\alph*)]
\item Show that two different sequences of signs give two different values of $X$.
\item Show that the law of $X$ has no atom.
\item Show that $X$ takes its values in a compact set $C$ that can be covered, for each $n$, by $2^n$ intervals of length at most $\frac2{n\cdot n!}$.
\item Deduce that $C$ has Hausdorff dimension $0$, and that the law of $X$ is continuous and singular.
\end{enumerate}
\end{exercise}

\begin{exercise}[$\star\star$]\label{pb:smooth}
\emph{A smooth sum.} Let $X=\sum_{n\ge1}\frac{\varepsilon_n}{n^2}$ and $\varphi(t)=\mathbb Ee^{itX}$.
\begin{enumerate}[label=(\alph*)]
\item Show that $\varphi(t)=\prod_{n\ge1}\cos\frac t{n^2}$.
\item Show that $|\varphi(t)|\le e^{-\sqrt{|t|}/48}$ for $|t|\ge1$.
\item Deduce that $X$ has a density which is infinitely differentiable.
\item More generally, show that $\sum_n\varepsilon_nn^{-\alpha}$ converges almost surely and has an infinitely differentiable density for every $\alpha>\frac12$.
Compare with $\sum_n\varepsilon_n2^{-n}$, which is uniform on $[-1,1]$.
\end{enumerate}
\end{exercise}

\begin{exercise}[$\star\star$]\label{pb:purity}
\emph{All or nothing.} Let $X_1,X_2,\dots$ be independent random variables, each taking countably many values, such that $X=\sum_nX_n$ converges almost surely.
Let $\mu$ be the law of $X$.
\begin{enumerate}[label=(\alph*)]
\item Show that $\mu$ is \emph{pure}: either it is carried by a countable set, or it is continuous and singular, or it is absolutely continuous.
\item Let $0<\lambda<1$ and $X_\lambda=\sum_{n\ge1}\varepsilon_n\lambda^n$. Show that the law of $X_\lambda$ has no atom. (Look at the atoms of largest mass.)
\item Deduce that the law of $X_\lambda$ is either absolutely continuous or singular, and decide which for $\lambda=\frac12$ and $\lambda=\frac13$.
\end{enumerate}
\end{exercise}

\begin{exercise}[$\star\star\star$]\label{pb:golden}
\emph{The golden ratio refuses a density.} Let $\varphi=\frac{1+\sqrt5}2$ be the golden ratio and $\mu$ the law of $Y=\sum_{n\ge1}\varepsilon_n\varphi^{-n}$, with
characteristic function $\hat\mu(t)=\mathbb Ee^{itY}$.
\begin{enumerate}[label=(\alph*)]
\item Show that $L_m=\varphi^m+\bigl(-\frac1\varphi\bigr)^m$ is an integer for every $m\ge0$.
\item Show that $|\hat\mu(\pi\varphi^k)|\to c^2$ as $k\to\infty$, where $c=\prod_{j\ge1}\bigl|\cos\frac\pi{\varphi^j}\bigr|>0$.
\item Deduce that $\mu$ is continuous and singular, although $\frac1\varphi>\frac12$ and the $2^n$ partial sums are spread over an interval of length about $2\varphi$.
\item Let $\theta>1$ be an algebraic integer which is not an integer, whose other conjugates all have modulus $<1$ (a \emph{Pisot number}). Show that the law of
$\sum_n\varepsilon_n\theta^{-n}$ is singular too.
\end{enumerate}
\end{exercise}

\begin{exercise}[Journey]\label{pb:kaczeros}
\emph{How many real roots?} Let $a_0,\dots,a_n$ be independent standard Gaussian variables, $f(t)=\sum_{j=0}^na_jt^j$, and $N_n$ the number of real roots of $f$.
\begin{enumerate}[label=(\alph*)]
\item \emph{Kac's counting formula.} Let $g$ be continuously differentiable on $[a,b]$, with $g(a)g(b)\ne0$ and $g'(t)\ne0$ whenever $g(t)=0$. Show that the number
of zeros of $g$ in $[a,b]$ is $\lim_{\varepsilon\to0}\frac1{2\varepsilon}\int_a^b\mathbf 1_{|g(t)|<\varepsilon}|g'(t)|\dd t$.
\item Show that, almost surely, $f$ has no multiple root, and that $\frac1{2\varepsilon}\int_a^b\mathbf 1_{|f(t)|<\varepsilon}|f'(t)|\dd t\le n$ for every $\varepsilon>0$.
\item Let $A(t)=\sum_jt^{2j}$, $B(t)=\sum_jjt^{2j-1}$, $C(t)=\sum_jj^2t^{2j-2}$. Show that the number of real roots in $[a,b]$ has mean
\[
\int_a^b\rho_n(t)\dd t,\qquad \rho_n=\frac{\sqrt{AC-B^2}}{\pi A}.
\]
\item Show that $\pi^2\rho_n(t)^2=\frac1{4t}\frac{\dd}{\dd t}\Bigl(t\frac{A'(t)}{A(t)}\Bigr)$, and deduce \emph{Kac's density}
\[
\rho_n(t)=\frac1\pi\sqrt{\frac1{(1-t^2)^2}-\frac{(n+1)^2t^{2n}}{(1-t^{2n+2})^2}}.
\]
\item Show that the numbers of roots in $]-\infty,-1[$, $]-1,0[$, $]0,1[$, $]1,\infty[$ have the same law, and that $\mathbb EN_n\sim\frac2\pi\log n$.
\item A random polynomial of degree $100$ thus has about three real roots, while a polynomial of degree $100$ can have all its roots real. Where are the
other roots?
\end{enumerate}
\end{exercise}

\section*{Hints}

\begin{hint}{pb:pairwisebits}
For (b), the sets whose index lies in a block $[a,a+2^m[$ with $2^m\mid a$ are $S'\cup U$ with $S'$ fixed and $U\subseteq\{1,\dots,m\}$ free. Factor the block sum.
\end{hint}

\begin{hint}{pb:sawtooth}
For (a), expand $B$ in Fourier series. For (b), $B(2^kx)$ depends on all the binary digits of $x$ after the $k$-th, but only slightly on the far ones:
replace it by a function of the next $L$ digits and control the error in $L^2$.
\end{hint}

\begin{hint}{pb:peaks}
The peak indicators at $i$ and $i+1$ are never both $1$; those at $i$ and $i+2$ are both $1$ for $16$ of the $120$ orders of five numbers. For (b), cut
$\{1,\dots,n\}$ into long blocks separated by short gaps.
\end{hint}

\begin{hint}{pb:tilings}
For (a), a tiling with $s$ squares and $d$ dominoes has probability $\varphi^{-s-2d}$. For (b), $u_n=P(H_n)$ satisfies a linear recurrence. For (c), first prove
the central limit theorem for the number of dominoes among the tiles ending before $n$, without conditioning.
\end{hint}

\begin{hint}{pb:harper}
For (a), factor $G$. For (b), $x(A+A')=xe^{-x}(e^xA)'$ and Rolle's theorem. For (c), differentiate the recurrence once and twice at $x=1$.
For (d), the ratio $P(M=m+1)/P(M=m)$ is decreasing and close to $1$ for $m$ near the mode.
\end{hint}

\begin{hint}{pb:fixedtv}
Write both laws with the alternating series of $e^{-1}$, and note the sign of each difference.
\end{hint}

\begin{hint}{pb:kubslow}
For (c), the law of $W_x$ sits on a lattice of step $\frac1{\sigma_x}$, and some atom is large.
\end{hint}

\begin{hint}{pb:twelve}
For (a), replace the $X_i$ by Gaussian variables one at a time and expand $h$ to order four. For (c), $P(U_1+\dots+U_{12}<2)$ is the volume of a corner of the cube.
\end{hint}

\begin{hint}{pb:threeap}
Apply the Chen--Stein theorem with the progressions meeting a given one as its neighbourhood. How many progressions contain a given pair?
\end{hint}

\begin{hint}{pb:spearman}
For (c), compute $\mathbb E[D^2\mid\pi]$ exactly: it is $\frac{2\lambda}4\bigl(2\sum_ib_i^2c_{\pi(i)}^2+\dots\bigr)$, and its fluctuations are those of another statistic of the
same kind as $W$.
\end{hint}

\begin{hint}{pb:erdostaylor}
For (a), rotate the lattice by $45^\circ$. For (b), expand $\mathbb E\binom{N_n}k$ with the Markov property. For (c), use the factorial moments to compute
$\mathbb E(1-\frac s{G_n})^{N_n}$.
\end{hint}

\begin{hint}{pb:dice}
Use the characters of $\Z/7\Z$: what is $\sum_{f=1}^6\omega^{f}$ for a seventh root of unity $\omega\ne1$?
\end{hint}

\begin{hint}{pb:squaresmodp}
Use the upper bound lemma. For $|\sum_a e(ha^2/p)|^2$, expand the square and substitute $a=b+c$.
\end{hint}

\begin{hint}{pb:lucky}
$L_m$ is the number of pairs of $m$-digit strings with the same digit sum. For (c), expand $\log\mathbb Ee^{itD}$ to order $t^4$.
\end{hint}

\begin{hint}{pb:signedsums}
The integrand has two peaks, at $t=0$ and $t=\pm\pi$.
\end{hint}

\begin{hint}{pb:smirnov}
For (a), reflect a path at its successive visits to $h$, $-h$, $h$, \dots, and use inclusion and exclusion. For (d), Poisson summation.
\end{hint}

\begin{hint}{pb:primeseries}
Kolmogorov's three series, and Mertens's theorem on dyadic blocks $e^{2^j}<p\le e^{2^{j+1}}$.
\end{hint}

\begin{hint}{pb:factorialseries}
$\sum_{k>n}\frac1{k!}<\frac1{n\cdot n!}$.
\end{hint}

\begin{hint}{pb:smooth}
$\log\cos u\le-\frac{u^2}2$ for $|u|<\frac\pi2$. Only the factors with $n\ge\sqrt{|t|}$ are needed.
\end{hint}

\begin{hint}{pb:purity}
Let $H$ be the group generated by all the values of all the $X_n$. Is the event $\{X\in E+H\}$ a tail event?
\end{hint}

\begin{hint}{pb:golden}
Split $\hat\mu(\pi\varphi^k)$ as $\prod_{m=0}^{k-1}\cos(\pi\varphi^m)$ times $\prod_{j\ge1}\cos(\pi\varphi^{-j})$, and note that $\pi\varphi^m$ is close to a multiple of $\pi$.
\end{hint}

\begin{hint}{pb:kaczeros}
For (b), between two zeros of $f'$ the function $f$ is monotone. For (c), condition on $f(t)$. For (e), reverse the coefficients.
\end{hint}

\section*{Solutions}

\begin{solution}{pb:pairwisebits}
\idea Products of signs are pairwise independent but very far from independent: the sum of all of them over $k$ signs is a product, $\prod(1+\varepsilon_i)-1$,
which is either huge or $-1$. Variance bounds survive pairwise independence; the Gaussian shape does not.

\writeup (a) If $i\in S$, changing the sign of $\varepsilon_i$ is a bijection of the outcomes that preserves probabilities and changes the sign of $\chi_S$, so
$P(\chi_S=1)=\frac12$. For $S\ne T$, $\chi_S\chi_T=\chi_{S\triangle T}$ because $\varepsilon_i^2=1$, and $S\triangle T\ne\emptyset$, so $\mathbb E\chi_S\chi_T=0$. Two signs of
mean $0$ that are uncorrelated are independent: $P(\chi_S=1,\chi_T=1)=\frac14\mathbb E(1+\chi_S)(1+\chi_T)=\frac14$, and similarly for the other three cases. But
$\chi_{\{1,2\}}=\chi_{\{1\}}\chi_{\{2\}}$ is a function of the other two, so the three are not independent.

(b) Write $\operatorname{ind}(S)=\sum_{i\in S}2^{i-1}$, so that $\operatorname{ind}$ is a bijection from finite sets to $\N$, with $\operatorname{ind}(\emptyset)=0$, and
$T_N+1=\sum_{\operatorname{ind}(S)\le N}\chi_S$ (the empty set contributes $\chi_\emptyset=1$). For $N=2^k-1$ the sets of index $\le N$ are the subsets of $\{1,\dots,k\}$, and
$\sum_{S\subseteq\{1,\dots,k\}}\chi_S=\prod_{i\le k}(1+\varepsilon_i)$, which is $2^k$ if $\varepsilon_1=\dots=\varepsilon_k=1$ and $0$ otherwise. So $T_N\in\{2^k-1,-1\}$.

For general $N$, write $N+1=2^{m_1}+\dots+2^{m_r}$ with $m_1>\dots>m_r$. Then $\{0,\dots,N\}$ is the union of the blocks $[a_j,a_j+2^{m_j}[$ with
$a_j=2^{m_1}+\dots+2^{m_{j-1}}$. Since $2^{m_j}$ divides $a_j$, the sets with index in the $j$-th block are $S_j\cup U$, where $S_j$ is the set of index $a_j$ (contained in
$\{m_j+1,m_j+2,\dots\}$) and $U$ runs over the subsets of $\{1,\dots,m_j\}$. So the block contributes $\chi_{S_j}\prod_{i\le m_j}(1+\varepsilon_i)$. Let $i_0$ be the first $i$
with $\varepsilon_i=-1$; it is finite almost surely. A block with $m_j\ge i_0$ contributes $0$; the other blocks have distinct $m_j<i_0$, so there are at most $i_0$ of them,
each contributing at most $2^{i_0}$ in absolute value. Hence $|T_N+1|\le i_0\,2^{i_0}$ for every $N$: almost surely $(T_N)$ is bounded, and $\frac{T_N}N\to0$.

(c) Along $N=2^k-1$, $P\bigl(T_N=2^k-1\bigr)=2^{-k}\to0$ and otherwise $\frac{T_N}{\sqrt N}=-\frac1{\sqrt N}\to0$. So $\frac{T_N}{\sqrt N}\to0$ in probability: the limit is the
point mass at $0$, not $N(0,1)$. (By (b) this even holds almost surely along all $N$.)

(d) Take $U_1,\dots,U_k$ independent and uniform on $[0,1[$, and for $\emptyset\ne S\subseteq\{1,\dots,k\}$ put $V_S=\bigl\{\sum_{i\in S}U_i\bigr\}$. Each $V_S$ is uniform.
If $S\ne T$, pick $i$ in one of them and not in the other, say $i\in S\setminus T$: given all $U_j$, $j\ne i$, the variable $V_S$ is a rotation of the uniform variable
$U_i$, hence uniform, while $V_T$ is fixed. So $V_S$ is uniform and independent of $V_T$: the $N=2^k-1$ points $V_S$ are pairwise independent and uniform. For the average
$A_N$ of the $f(V_S)$, pairwise independence is all that the variance needs: $\operatorname{Var}A_N=\frac1N\operatorname{Var}f(U)$, and Chebyshev's inequality gives the
same guarantee as $N$ independent samples, from $k=\log_2(N+1)$ random numbers. A Gaussian confidence interval, on the other hand, rests on the central limit theorem
for $A_N$, and (c) is the same construction in the group $\{\pm1\}$ instead of the circle: there the error is $-\frac1N$ with probability $1-2^{-k}$ and huge otherwise,
nothing like a bell curve.

\tools{Volume~4, Chapter~5 (independence and variance of sums), Chebyshev's inequality.}
\leads Pairwise independent hash functions built this way are a basic tool of randomised algorithms: they save randomness wherever only second moments
matter.
\end{solution}

\begin{solution}{pb:sawtooth}
\idea The doubling map $x\mapsto2x$ shifts the binary digits, so $B(2^kx)$ depends on the digits after the $k$-th, mostly on the next few. The terms are
\emph{almost} independent, the central limit theorem holds, but the correlations do not vanish and they change the variance: $\frac14$ instead of $\frac1{12}$.

\writeup (a) $B(x)=-\sum_{h\ge1}\frac{\sin2\pi hx}{\pi h}$ in $L^2(0,1)$, and $B(mx)=-\sum_{h'\ge1}\frac{\sin2\pi mh'x}{\pi h'}$. By Parseval's formula (Volume~5, Theorem~8.11), only the
frequencies $h=mh'$ survive, each contributing $\frac12\cdot\frac1{\pi h}\cdot\frac1{\pi h'}$:
\[
\int_0^1B(x)B(mx)\dd x=\sum_{h'\ge1}\frac1{2\pi^2mh'^2}=\frac1{2\pi^2m}\cdot\frac{\pi^2}6=\frac1{12m}.
\]

(b) \emph{Variance.} Since $x\mapsto\{2^jx\}$ preserves Lebesgue measure and $B$ has period $1$, $\int B(2^jx)B(2^kx)\dd x=\int B(y)B(2^{k-j}y)\dd y=\frac{2^{-(k-j)}}{12}$ for
$k\ge j$. With $S_n=\sum_{k<n}B(2^kx)$,
\[
\frac{\operatorname{Var}S_n}n=\frac1{12}+\frac2n\sum_{d=1}^{n-1}(n-d)\frac{2^{-d}}{12}\longrightarrow\frac1{12}+\frac2{12}\sum_{d\ge1}2^{-d}=\frac14=\sigma^2.
\]

\emph{Truncation.} Let $d_1,d_2,\dots$ be the binary digits of $x$, independent fair bits. Then $\{2^kx\}=\sum_{i\ge1}d_{k+i}2^{-i}$. Put
$B_L(2^kx)=\sum_{i=1}^Ld_{k+i}2^{-i}+2^{-L-1}-\frac12$, a centred function of $d_{k+1},\dots,d_{k+L}$. The remainder is
$B(2^kx)-B_L(2^kx)=2^{-L}\bigl(\{2^{k+L}x\}-\frac12\bigr)=2^{-L}B(2^{k+L}x)$, so with $W_n=\frac{S_n}{\sqrt n}$ and $W_n^L=\frac1{\sqrt n}\sum_{k<n}B_L(2^kx)$, the variance
computation (shifted by $L$) gives $\|W_n-W_n^L\|_2=2^{-L}\sigma_n$, where $\sigma_n^2=\frac{\operatorname{Var}S_n}n\to\frac14$; and $v_{n,L}=\operatorname{Var}W_n^L$ satisfies
$|\sqrt{v_{n,L}}-\sigma_n|\le2^{-L}\sigma_n$.

\emph{The limit.} The variables $B_L(2^kx)$ are bounded and $(L-1)$-dependent, and $v_{n,L}\ge\frac1{16}$ for $L\ge2$ and $n$ large, so by Problem~\ref{pb:peaks}(b),
$W_n^L/\sqrt{v_{n,L}}\Rightarrow N(0,1)$. Characteristic functions then converge uniformly on compact sets, so $\mathbb Ee^{itW_n^L}-e^{-t^2v_{n,L}/2}\to0$ for each $t$. Hence
\[
\bigl|\mathbb Ee^{itW_n}-e^{-t^2/8}\bigr|\le|t|\,\|W_n-W_n^L\|_2+\bigl|\mathbb Ee^{itW_n^L}-e^{-t^2v_{n,L}/2}\bigr|+\frac{t^2}2\bigl|v_{n,L}-\tfrac14\bigr|,
\]
whose $\limsup$ as $n\to\infty$ is at most $(|t|+3t^2)2^{-L}$. As $L$ is arbitrary, $\mathbb Ee^{itW_n}\to e^{-t^2/8}$, and $W_n\Rightarrow N(0,\frac14)$ by Lévy's continuity
theorem (Theorem~3.3).

(c) The same proof in base $3$ ($B(3^kx)$ depends on the ternary digits after the $k$-th) gives the limit $N(0,\sigma^2)$ with
$\sigma^2=\frac1{12}+\frac2{12}\sum_{d\ge1}3^{-d}=\frac1{12}+\frac1{12}=\frac16$.

(d) The variance computation of (b) is exactly $\sigma^2=\int B^2+2\sum_{d\ge1}\int B(x)B(2^dx)\dd x$. Take $f(x)=g(x)-g(2x)$ with $g(x)=\cos2\pi x$, a trigonometric
polynomial of mean $0$. Then $\sum_{k<n}f(2^kx)=g(x)-g(2^nx)$ telescopes, is bounded by $2$, and $\frac1{\sqrt n}\sum_{k<n}f(2^kx)\to0$ surely. Here the correlations
cancel the variance completely.

\tools{Volume~5, Theorem~8.11 (Parseval); Theorem~3.3 (Lévy's continuity theorem); Problem~\ref{pb:peaks}(b).}
\leads This is Kac's theorem of 1946 for lacunary sums $\sum f(2^kx)$: the limit is Gaussian with variance $\int f^2+2\sum_d\int f(x)f(2^dx)\dd x$, and the variance is $0$
exactly for coboundaries $g(x)-g(2x)$, as in (d). For a non-lacunary sequence such as $f(kx)$ nothing of the sort holds.
\end{solution}

\begin{solution}{pb:peaks}
\idea Peaks are local events: the indicator of a peak at $i$ depends on three consecutive values, so peaks two apart or more are independent. Count the
correlations of neighbours, then cut the line into long blocks separated by short gaps, so that the blocks are independent.

\writeup (a) Realise the permutation by independent uniform $U_1,\dots,U_n$ (the ranks of the $U_i$ form a uniform permutation), and let $X_i=\mathbf 1\{U_{i-1}<U_i>U_{i+1}\}$
for $1<i<n$. Each $X_i$ has mean $\frac13$ (the middle one of three is the largest with probability $\frac13$), so $\mathbb EP_n=\frac{n-2}3$. Adjacent peaks are
impossible, so $\operatorname{Cov}(X_i,X_{i+1})=-\frac19$; peaks at $i$ and $i+2$ mean $U_{i-1}<U_i>U_{i+1}<U_{i+2}>U_{i+3}$, which holds for $16$ of the $5!$ orders. (The largest of the five values sits in position $2$ or $4$. If it is in position $2$, the first value is any of
the other four, and of the last three the middle one must be the largest: $4\cdot2=8$ orders; position $4$ gives $8$ more by symmetry.) so $\operatorname{Cov}(X_i,X_{i+2})=\frac{16}{120}-\frac19=\frac1{45}$; and $X_i,X_j$ are independent for $|i-j|\ge3$. Hence, for $n\ge4$,
\[
\operatorname{Var}P_n=(n-2)\frac29-2(n-3)\frac19+2(n-4)\frac1{45}=\frac{10+2n-8}{45}=\frac{2(n+1)}{45}.
\]

(b) Let $r=\lfloor n^{1/4}\rfloor$ and cut $\{1,\dots,n\}$ into consecutive pieces: a \emph{big} block of length $r$, a \emph{small} block of length $m$, a big block, and so
on, plus a final piece of length less than $r+m$. Let $U_1,U_2,\dots$ be the sums over the big blocks and $V$ the sum of everything else. Two big blocks are separated by
at least $m$ positions, so the $U_j$ are independent; the same holds for the small blocks. There are at most $\frac n r+1$ small blocks, each of sum bounded by $mc$,
so $\operatorname{Var}V\le2\bigl(\frac nr+1\bigr)(mc)^2+2(r+m)^2c^2=o(n)$, using independence of the small block sums and treating the final piece separately. Hence
$\|S_n-\sum_jU_j\|_2=o(\sqrt n)=o(\sqrt{\operatorname{Var}S_n})$, and $\operatorname{Var}\sum_jU_j=\operatorname{Var}S_n\,(1+o(1))$. Each $|U_j|\le rc=o(\sqrt{\operatorname{Var}S_n})$,
so for every $\varepsilon>0$ the Lindeberg sum $\sum_j\mathbb E[U_j^2;|U_j|>\varepsilon\sqrt{\operatorname{Var}S_n}]$ is $0$ for $n$ large: by Lindeberg's theorem (Volume~4,
Theorem~6.10), $\sum_jU_j/\sqrt{\operatorname{Var}S_n}\Rightarrow N(0,1)$, and so does $S_n/\sqrt{\operatorname{Var}S_n}$, since the difference tends to $0$ in $L^2$.

(c) The $X_i-\frac13$ are centred, bounded and $2$-dependent, with $\operatorname{Var}P_n=\frac{2(n+1)}{45}\ge\frac{2n}{45}$. By (b), $\bigl(P_n-\frac{n-2}3\bigr)/\sqrt{\operatorname{Var}P_n}\Rightarrow N(0,1)$,
and $\operatorname{Var}P_n\sim\frac{2n}{45}$.

\tools{Volume~4, Theorem~6.10 (Lindeberg); Volume~4, Chapter~5 (variance of sums).}
\leads The same $m$-dependent central limit theorem gives Gaussian laws for the number of descents, of runs, and of occurrences of a pattern in a random word. Rates come from Stein's method (Volume~11, Chapter~6), which handles local dependence directly.
\end{solution}

\begin{solution}{pb:tilings}
\idea A uniform tiling is a sequence of independent tiles conditioned to stop exactly at $n$, and the conditioning costs almost nothing: the event has probability
close to $\frac1{1+q}$ whatever happened long before. So the dominoes behave like a renewal--reward process, whose fluctuations are Gaussian.

\writeup (a) A tiling of $[0,n]$ with $s$ squares and $d$ dominoes ($s+2d=n$) is produced by the first $s+d$ tiles with probability
$(1-q)^sq^d=\varphi^{-s}\varphi^{-2d}=\varphi^{-n}$ (note $1-q=1-\varphi^{-2}=\varphi^{-1}$, since $\varphi^2=\varphi+1$). This does not depend on the tiling, and the
events for different tilings are disjoint with union $H_n$. So given $H_n$ the tiling is uniform, and $P(H_n)=F_{n+1}\varphi^{-n}$.

(b) $u_n=P(H_n)$ satisfies $u_0=1$, $u_1=1-q$, and $u_n=(1-q)u_{n-1}+qu_{n-2}$ by conditioning on the last tile ending at $n$. The characteristic equation
$x^2=(1-q)x+q$ has roots $1$ and $-q$, so $u_n=\alpha+\beta(-q)^n$, and the initial values give $\alpha=\frac1{1+q}$, $\beta=\frac q{1+q}$: $u_n=\frac{1-(-q)^{n+1}}{1+q}$.

(c) \emph{Without conditioning.} Let $\tau_j=T_1+\dots+T_j$, $\nu(m)=\max\{j:\tau_j\le m\}$, and $D'_m$ the number of dominoes among $T_1,\dots,T_{\nu(m)}$. With
$a=\frac q{1+q}$ and $Y_j=\mathbf 1\{T_j=2\}-aT_j$, we have $\mathbb EY_j=q-a(1+q)=0$ and
\[
D'_m-am=\sum_{j\le\nu(m)}Y_j+a\bigl(\tau_{\nu(m)}-m\bigr),\qquad\bigl|a(\tau_{\nu(m)}-m)\bigr|\le a.
\]
Let $\mu=\mathbb ET=1+q$ and $k_m=\lfloor m/\mu\rfloor$. By the strong law of renewal (Theorem~2.7), $\nu(m)/m\to\frac1\mu$, so $P(|\nu(m)-k_m|>\varepsilon m)\to0$
for every $\varepsilon>0$. On the complement, $\bigl|\sum_{j\le\nu(m)}Y_j-\sum_{j\le k_m}Y_j\bigr|$ is at most the maximum of the partial sums of $Y_{k_m+1},Y_{k_m+2},\dots$
(or of $Y_{k_m},Y_{k_m-1},\dots$) over $\varepsilon m$ terms, and Kolmogorov's inequality (Volume~4, Corollary~10.16) bounds the probability that it exceeds $\eta\sqrt m$ by
$2\varepsilon\operatorname{Var}Y/\eta^2$. With the central limit theorem for $\sum_{j\le k_m}Y_j$, this gives $\frac{D'_m-am}{\sqrt m}\Rightarrow N(0,s^2)$ with
$s^2=\frac{\operatorname{Var}Y}\mu$. Now $Y=(1-a)\mathbf 1\{T=2\}-a$, so $\operatorname{Var}Y=(1-a)^2q(1-q)$ and $s^2=\frac{q(1-q)}{(1+q)^3}$. Since
$1+q=\frac{\varphi^2+1}{\varphi^2}=\frac{\sqrt5\,\varphi}{\varphi^2}=\frac{\sqrt5}\varphi$ (because $\varphi^2+1=\varphi+2=\sqrt5\,\varphi$), we get
$s^2=\frac{\varphi^{-3}}{5\sqrt5\,\varphi^{-3}}=\frac1{5\sqrt5}$ and $a=\frac{\varphi^{-2}\varphi}{\sqrt5}=\frac1{\sqrt5\,\varphi}=\frac{5-\sqrt5}{10}$.

\emph{With conditioning.} Let $L=\lfloor n^{1/4}\rfloor$ and $J$ the first $j$ with $\tau_j\ge n-L$, so $\tau_J\in\{n-L,n-L+1\}$. The variable $D'_{n-L}$ is a function of
$T_1,\dots,T_J$, the tiles after $J$ are independent of them, and so $P(H_n\mid T_1,\dots,T_J)=u_{n-\tau_J}=\frac1{1+q}+O(q^{L-1})$. Hence
$P(D'_{n-L}\le y\mid H_n)=P(D'_{n-L}\le y)+O(q^{L-1})$ uniformly in $y$, because $u_n=\frac1{1+q}+O(q^{n+1})$ too. On $H_n$, $D_n$ is $D'_n$, which differs from $D'_{n-L}$ by
at most $L=o(\sqrt n)$. So $\frac{D_n-an}{\sqrt n}\Rightarrow N(0,s^2)$.

\tools{Theorem~2.7 (renewal); Volume~4, Corollary~10.16 (Kolmogorov's inequality) and Theorem~6.10 (central limit theorem).}
\leads In combinatorial language, $P(D_n=k)=\binom{n-k}k/F_{n+1}$, and the result says that the coefficients of the rows of Pascal's triangle read along the
diagonals that give the Fibonacci numbers are Gaussian. The probabilistic proof gives the same answer for tilings with any finite set of tile lengths.
\end{solution}

\begin{solution}{pb:harper}
\idea A probability generating function with real roots factors into linear factors, and each factor is a coin. So a law with real-rooted generating function is a sum
of independent coins, and the central limit theorem applies as soon as its variance tends to infinity. Dobiński's representation computes the variance.

\writeup (a) $G$ has non-negative coefficients, so it has no positive root, and $G(0)=P(X=0)$. Its roots are $-r_1,\dots,-r_d$ with $r_i\ge0$, $d\le n$, and
$G(x)=c\prod_i(x+r_i)$. Since $G(1)=1$, $G(x)=\prod_i\frac{x+r_i}{1+r_i}=\prod_i(1-p_i+p_ix)$ with $p_i=\frac1{1+r_i}\in\left]0,1\right]$, the generating function of a sum of
independent Bernoulli$(p_i)$ variables; add $n-d$ Bernoulli variables equal to $0$.

(b) $S(n+1,k)=kS(n,k)+S(n,k-1)$ (the element $n+1$ joins one of the $k$ blocks of a partition of $\{1,\dots,n\}$, or forms its own block), which is
$A_{n+1}=xA_n'+xA_n$. Now $A_n+A_n'=e^{-x}(e^xA_n)'$. Suppose $A_n$ has $n$ distinct real roots $r_n<\dots<r_1=0$ ($A_n(0)=S(n,0)=0$ for $n\ge1$). Then $e^xA_n$ vanishes at the
$r_i$ and tends to $0$ at $-\infty$; by Rolle's theorem $(e^xA_n)'$ has a root in each $]r_{i+1},r_i[$ and one in $]-\infty,r_n[$: $n$ distinct roots, all negative, and these are
all the roots of $A_n+A_n'$, which has degree $n$. Multiplying by $x$, $A_{n+1}$ has the $n+1$ distinct roots $0$ and these. Start from $A_1=x$.

(c) At $x=1$ the recurrence gives $B_{n+1}=B_n+A_n'(1)$, so $\mathbb EK_n=\frac{A_n'(1)}{B_n}=\frac{B_{n+1}}{B_n}-1$. Differentiating, $A_{n+1}'=A_n+A_n'+x(A_n'+A_n'')$, so
$A_{n+1}'(1)=B_n+2A_n'(1)+A_n''(1)$, and $A_{n+1}'(1)=B_{n+2}-B_{n+1}$; hence $A_n''(1)=B_{n+2}-3B_{n+1}+B_n$ and
$\mathbb EK_n^2=\frac{A_n''(1)+A_n'(1)}{B_n}=\frac{B_{n+2}-2B_{n+1}}{B_n}$. On the other side, Dobiński's formula $B_k=\frac1e\sum_m\frac{m^k}{m!}$ gives
$\mathbb EM^j=\frac{B_{n+j}}{B_n}$. So $\mathbb EK_n=\mathbb EM-1$, and
\[
\operatorname{Var}K_n=\frac{B_{n+2}-2B_{n+1}}{B_n}-\Bigl(\frac{B_{n+1}}{B_n}-1\Bigr)^2=\frac{B_{n+2}}{B_n}-\Bigl(\frac{B_{n+1}}{B_n}\Bigr)^2-1=\operatorname{Var}M-1.
\]

(d) Let $\pi_m=P(M=m)$ and $r(m)=\frac{\pi_{m+1}}{\pi_m}=\frac{(1+1/m)^n}{m+1}$, which decreases in $m$. Let $m_0$ be a mode: $r(m_0-1)\ge1\ge r(m_0)$. From
$e^{n/(m_0+1)}\le(1+\frac1{m_0})^n\le m_0+1$ and $(1+\frac1{m_0-1})^n\ge m_0$ we get $(m_0-1)\log m_0\le n\le(m_0+1)\log(m_0+1)$, so $m_0\sim\frac n{\log n}$.
For $m\ge2$, $\log r(m-1)-\log r(m)=n\log\bigl(1+\frac1{m^2-1}\bigr)+\log\frac{m+1}m\le\frac{2n}{m^2}+\frac1m$. Let $\varepsilon=\frac{8n}{m_0^2}+\frac2{m_0}$, which tends to $0$,
and $J=\lfloor\varepsilon^{-1/2}\rfloor\le\frac{m_0}2$ for $n$ large. For $|i|\le J$ the steps between $m_0$ and $m_0+i$ have $|\log r|\le(J+1)\varepsilon$ (start from
$\log r(m_0)\le0\le\log r(m_0-1)$, with consecutive differences at most $\varepsilon$), so $\pi_{m_0+i}\ge\pi_{m_0}e^{-J(J+1)\varepsilon}\ge e^{-2}\pi_{m_0}$, while $\pi_m\le\pi_{m_0}$ for
all $m$. Now fix any real $c$. At most $\frac J2+1$ integers lie in $]c-\frac J4,c+\frac J4[$, each of probability at most $\pi_{m_0}$, while at least $J$ integers of $[m_0-J,m_0+J]$ lie
outside, each of probability at least $e^{-2}\pi_{m_0}$. So $P(|M-c|\ge\frac J4)\ge e^{-2}P(|M-c|<\frac J4)$ for $J\ge2$, hence $P(|M-c|\ge\frac J4)\ge\frac{e^{-2}}{1+e^{-2}}$, and
$\mathbb E(M-c)^2\ge\frac{J^2}{16}\cdot\frac{e^{-2}}{1+e^{-2}}$. Taking $c=\mathbb EM$: $\operatorname{Var}M\to\infty$, since $J\to\infty$.

So $\operatorname{Var}K_n=\sum_ip_i(1-p_i)\to\infty$ for the Bernoulli decomposition of (a), (b). The summands are bounded by $1$, so Lindeberg's condition holds trivially
(Volume~4, Theorem~6.10), and $\frac{K_n-\mathbb EK_n}{\sqrt{\operatorname{Var}K_n}}\Rightarrow N(0,1)$.

\tools{Volume~1, Theorem~3.10 (Rolle); Volume~5, Exercise~2.3 (Dobiński's formula); Volume~4, Theorem~6.10 (Lindeberg).}
\leads Harper proved this in 1967, and the method is now a standard trick: matching polynomials of graphs, independence polynomials of claw-free graphs, Eulerian
polynomials all have real roots, and all give Gaussian laws for free. Real roots give more: Darroch's theorem says that the mode of a sum of independent coins is within
$1$ of the mean (a direct computation verifies it for $K_n$).
\end{solution}

\begin{solution}{pb:fixedtv}
\idea Inclusion and exclusion gives the law of $F_n$ as a truncated Poisson series: the two laws differ only by the tails of the series for $e^{-1}$, which are tiny. A
general-purpose method cannot compete with an exact formula.

\writeup (a) By inclusion and exclusion (Volume~4, Theorem~2.13), $P(F_n=k)=\frac1{k!}\sum_{j=0}^{n-k}\frac{(-1)^j}{j!}$ for $0\le k\le n$, while $e^{-1}/k!=\frac1{k!}\sum_{j\ge0}\frac{(-1)^j}{j!}$.
So $d_k=P(F_n=k)-\frac{e^{-1}}{k!}=-\frac1{k!}\sum_{j>n-k}\frac{(-1)^j}{j!}$, an alternating tail with decreasing terms: $d_k$ has the sign of $(-1)^{n-k}$ and
$|d_k|\le\frac1{k!\,(n-k+1)!}$. As both laws have total mass $1$, the distance is the sum of the positive parts of the differences, and the Poisson masses above $n$ only
count negatively. So
\[
\begin{aligned}
\|\mathcal L(F_n)-\text{Poisson}(1)\|&=\sum_{\substack{k\le n\\ n-k\text{ even}}}d_k\le\sum_{\substack{k\le n\\ n-k\text{ even}}}\frac1{k!\,(n+1-k)!}\\
&=\frac1{(n+1)!}\sum_{\substack{1\le j\le n+1\\ j\text{ odd}}}\binom{n+1}j=\frac{2^n}{(n+1)!}.
\end{aligned}
\]

(b) The indicator $I_i$ of $\{\pi(i)=i\}$ is not independent of $(I_j)_{j\ne i}$ (if all other points are fixed, so is $i$), so the neighbourhood of each $i$ in Volume~11,
Theorem~6.8, must be everything. Then $b_1=\bigl(\sum_ip_i\bigr)^2=1$, $b_2=\sum_{i\ne j}P(I_iI_j=1)=n(n-1)\frac1{n(n-1)}=1$, and the bound is $2(1-e^{-1})\approx1.26$:
useless. Local dependence is the wrong picture here: all the indicators are weakly dependent, none independent.

\tools{Volume~4, Theorem~2.13 (fixed points and inclusion--exclusion); Theorem~6.8 (Chen--Stein).}
\leads Coupling versions of the Chen--Stein method handle weak global dependence and give a bound of order $\frac1n$ for fixed points (Barbour, Holst and
Janson, Chapter~4), still far from the truth, which is superexponentially small.
\end{solution}

\begin{solution}{pb:kubslow}
\idea Stein's bound for the Wasserstein distance is of order $\frac1{\sigma_x}$, which is $(\log\log x)^{-1/2}$: correct, but $\log\log x$ grows so slowly that the
bound says nothing in practice. Lattice effects show that nothing better than $(\log\log x)^{-1/2}$ is possible in Kolmogorov's distance either.

\writeup (a) Let $X_p=(B_p-\frac1p)/\sigma_x$, independent and centred, with $\sum_p\mathbb EX_p^2=1$. Then
$\mathbb E|B_p-\frac1p|^3=\frac1p(1-\frac1p)\bigl((1-\frac1p)^2+\frac1{p^2}\bigr)\le\frac1p(1-\frac1p)$, so $\sum_p\mathbb E|X_p|^3\le\frac{\sigma_x^2}{\sigma_x^3}=\frac1{\sigma_x}$ and
Theorem~6.4 gives $d_W\le\frac3{\sigma_x}$. Proposition~6.3 there gives $d_K\le(2/\pi)^{1/4}\sqrt3\,\sigma_x^{-1/2}$, and $(2/\pi)^{1/4}\sqrt3=1.547\ldots$, while
$\sigma_x^2=\log\log x+O(1)$ by Mertens's theorem (Volume~6, Theorem~1.5). So $d_K\le2(\log\log x)^{-1/4}$ for $x$ large.

(b) $2(\log\log x)^{-1/4}\le0.01$ requires $\log\log x\ge200^4=1.6\times10^9$, that is $x\ge\exp(\exp(1.6\times10^9))$, a number with more than $10^{690\,000\,000}$ digits.

(c) $K=\sum_pB_p$ is integer-valued, so $W_x$ lives on a lattice of step $\frac1{\sigma_x}$. If $w$ is an atom of mass $m$, the distribution function jumps by $m$ at $w$
while $\Phi$ is continuous, so $|F(w)-\Phi(w)|+|F(w^-)-\Phi(w)|\ge m$ and $d_K\ge\frac m2$. By Chebyshev's inequality $P(|K-\mathbb EK|<2\sigma_x)\ge\frac34$, and at most
$4\sigma_x+1$ integers lie in that interval, so some atom has mass at least $\frac3{4(4\sigma_x+1)}\ge\frac1{8\sigma_x}$ for $\sigma_x\ge1$. Hence $d_K\ge\frac1{16\sigma_x}\ge\frac c{\sqrt{\log\log x}}$.
For $x=10^6$, $\sigma_x^2=2.435$, the largest atom gives $d_K\ge0.123$, and the exact distance computed numerically is $0.150$.

\tools{Proposition~6.3 and Theorem~6.4 (Stein); Volume~6, Theorem~1.5 (Mertens).}
\leads The same lower bound holds for $\omega(n)$ itself: Rényi and Turán (1958) proved that the error in the Erdős--Kac theorem is exactly of order
$(\log\log x)^{-1/2}$. Better approximations use the Poisson law, or the mod-Poisson convergence of Volume~11, Chapter~11, instead of the Gaussian.
\end{solution}

\begin{solution}{pb:twelve}
\idea Lindeberg's swapping argument with a fourth-order Taylor expansion: for symmetric pieces the first three moments of $X_i$ and of a Gaussian agree, so each swap
costs $O(n^{-2})$ and the total error is $O(\frac1n)$, better than the $n^{-1/2}$ of Berry--Esseen. But this is for smooth test functions; the tails tell another story.

\writeup (a) Let $Z_1,\dots,Z_n$ be standard Gaussian, independent of the $X_i$, so that $Z=\frac{Z_1+\dots+Z_n}{\sqrt n}$ is standard Gaussian. Put
$R_i=\frac1{\sqrt n}(X_1+\dots+X_{i-1}+Z_{i+1}+\dots+Z_n)$; then $\mathbb Eh(W)-\mathbb Eh(Z)=\sum_i\bigl(\mathbb Eh(R_i+\frac{X_i}{\sqrt n})-\mathbb Eh(R_i+\frac{Z_i}{\sqrt n})\bigr)$. Taylor's formula gives
$h(R+y)=h(R)+h'(R)y+\frac12h''(R)y^2+\frac16h'''(R)y^3+\rho$ with $|\rho|\le\frac{\|h^{(4)}\|}{24}y^4$. Since $R_i$ is independent of $X_i$ and $Z_i$, and $X_i$, $Z_i$
have the same moments of order $1$, $2$ and $3$ ($0$, $1$, $0$, by symmetry), the terms of order up to $3$ cancel and
$\bigl|\mathbb Eh(R_i+\frac{X_i}{\sqrt n})-\mathbb Eh(R_i+\frac{Z_i}{\sqrt n})\bigr|\le\frac{\|h^{(4)}\|}{24n^2}(m_4+3)$. Sum over $i$.

(b) $W=\frac1{\sqrt{12}}\sum_{i\le12}X_i$ with $X_i=\sqrt{12}(U_i-\frac12)$, symmetric, $\mathbb EX_i^2=1$, $m_4=144\,\mathbb E(U-\frac12)^4=\frac{144}{80}=\frac95$. With $h=\cos$,
$\|h^{(4)}\|=1$ and $n=12$: $|\mathbb E\cos W-\mathbb E\cos Z|\le\frac{9/5+3}{24\cdot12}=\frac1{60}$, and $\mathbb E\cos Z=e^{-1/2}$. Exactly,
$\mathbb E\cos W=\operatorname{Re}\prod_i\mathbb Ee^{i(U_i-1/2)}=\bigl(\int_{-1/2}^{1/2}\cos u\dd u\bigr)^{12}=(2\sin\tfrac12)^{12}=0.603968$, against $e^{-1/2}=0.606531$: the error
$0.0026$ is six times smaller than the bound.

(c) $|W|\le6$ always. $P(W>4)=P(\sum U_i>10)=P(\sum U_i<2)$ by $U_i\mapsto1-U_i$. The simplex $\{u\ge0,\sum u_i<2\}$ has volume $\frac{2^{12}}{12!}$; the part outside
the cube is where some $u_i>1$, and two coordinates cannot both exceed $1$; translating $u_i$ by $1$, each such part has volume $\frac1{12!}$. So
$P(W>4)=\frac{2^{12}-12}{12!}$ and $P(|W|>4)=\frac{2(2^{12}-12)}{12!}=1.705\times10^{-5}$, against $P(|Z|>4)=6.33\times10^{-5}$. Bound (a) controls smooth averages, not
probabilities of rare events: an indicator has no fourth derivative, and smoothing it on a scale $\delta$ costs a factor $\delta^{-4}$. A simulation that needs
events beyond $4\sigma$ is four times too optimistic with this generator, and events beyond $6\sigma$ never happen at all.

\tools{Volume~4, Theorem~6.10 (Lindeberg's swapping argument); Volume~4, Chapter~5 (uniform sums).}
\leads The rate $\frac1n$ for smooth functions and symmetric summands is the first case of the Edgeworth expansion. Exact methods cost little more: the Box--Muller transform
(1958) turns two uniform numbers into two independent Gaussian ones.
\end{solution}

\begin{solution}{pb:threeap}
\idea Progressions in a sparse random set are rare and almost independent events: two of them can only interact if they overlap, and overlapping pairs are few.
This is the situation of the Chen--Stein theorem, with the triangles of a random graph as the model (Example~6.10).

\writeup (a) A progression is determined by its first term and its difference, and $(a,d)$ and $(a+2d,-d)$ give the same set; for $N\ge5$ prime nothing else
does (the middle term $b$ of $\{x,b,y\}$ satisfies $2b=x+y$, and two different middles would force $3d=0$). So there are $\frac{N(N-1)}2$ progressions. Two different
sets of three elements share at most two.

(b) Index by progressions $\alpha$, $I_\alpha=\mathbf 1\{\alpha\subseteq A\}$, $p_\alpha=p^3$, and let $B_\alpha$ be the progressions meeting $\alpha$; $I_\alpha$ is independent of
the progressions disjoint from $\alpha$. A given residue $x$ lies in $(N-1)$ progressions as an end and $\frac{N-1}2$ as the middle, so $|B_\alpha|\le\frac{9(N-1)}2$. A given pair
$\{x,y\}$ lies in exactly $3$ progressions ($x$, $y$ as the ends, or one of them in the middle), so at most $6$ progressions share two elements with $\alpha$. With
$M=\frac{N(N-1)}2$ and $\lambda=Mp^3$,
\[
b_1\le M\frac{9(N-1)}2p^6,\qquad b_2\le M\Bigl(6p^4+\frac{9(N-1)}2p^5\Bigr),
\]
since two progressions sharing two elements involve four residues, and sharing one, five. By Theorem~6.8,
$\|\mathcal L(W)-\text{Poisson}(\lambda)\|\le\frac{b_1+b_2}\lambda\le\frac{9(N-1)}2(p^3+p^2)+6p$. With $p=cN^{-2/3}$: $Np^3=\frac{c^3}N$, $Np^2=c^2N^{-1/3}$, $p=cN^{-2/3}$, so the
distance is $O(N^{-1/3})$.

(c) $P(W=0)=e^{-\lambda}+O(N^{-1/3})$ and $\lambda=c^3\frac{N(N-1)}{2N^2}\to\frac{c^3}2$. And $|A|$ is binomial with mean $Np=cN^{1/3}\to\infty$ and variance at most the mean,
so $|A|/(cN^{1/3})\to1$ in probability by Chebyshev's inequality. For $N=10^4$ the bound is $0.09$; the simulation in one finds numerically $0.65$ against $e^{-1/2}=0.61$ for $c=1$.

\tools{Theorem~6.8 (Chen--Stein).}
\leads The threshold $N^{-2/3}$ is where a random set starts to contain progressions; Roth's theorem says that a set of density $\delta>0$ always contains one, and the
modern random versions of Roth's and Szemerédi's theorems (Conlon and Gowers, Schacht, 2016) find the threshold for ``every dense subset of a random set contains a
progression'': it is $N^{-1/2}$.
\end{solution}

\begin{solution}{pb:spearman}
\idea A random transposition changes $W$ a little and pulls it back towards $0$ on average, with $\lambda=\frac2{n-1}$. The exchangeable pair theorem then needs the
fluctuations of $\mathbb E[D^2\mid\pi]$, which is again a linear rank statistic, and of $|D|^3$.

\writeup (a) For $i\ne j$, $(\pi(i),\pi(j))$ is uniform over ordered pairs of distinct positions. Let $\bar y$ be the mean and $S_y=\sum_k(y_k-\bar y)^2$. Then
$\operatorname{Var}y_{\pi(i)}=\frac{S_y}n$ and $\operatorname{Cov}(y_{\pi(i)},y_{\pi(j)})=-\frac{S_y}{n(n-1)}$ (the sum of all $y_{\pi(i)}$ is constant). Hence
\[
\operatorname{Var}\sum_ix_iy_{\pi(i)}=\frac{S_y}n\sum_ix_i^2-\frac{S_y}{n(n-1)}\sum_{i\ne j}x_ix_j=\frac{S_y}{n(n-1)}\Bigl(n\sum_ix_i^2-\bigl(\sum_ix_i\bigr)^2\Bigr)=\frac{S_xS_y}{n-1}.
\]
For $W$, $\mathbb EW=\sum_ib_i\bar c=0$ and $\operatorname{Var}W=\frac{\sum b_i^2\sum c_j^2}{n-1}=1$.

(b) $\pi'$ is again uniform and $\pi=\pi'\circ(I\,J)$, so $(\pi,\pi')$ and $(\pi',\pi)$ have the same law, and so do $(W,W')$ and $(W',W)$. Given $\pi$, write $c'_i=c_{\pi(i)}$;
then $D=W'-W=b_Ic'_J+b_Jc'_I-b_Ic'_I-b_Jc'_J=-(b_I-b_J)(c'_I-c'_J)$. Averaging over the $n(n-1)$ ordered pairs, and using $\sum_{i\ne j}b_ic'_j=\bigl(\sum b\bigr)\bigl(\sum c'\bigr)-W=-W$:
\[
\mathbb E[D\mid\pi]=\frac{2(-W)-2(n-1)W}{n(n-1)}=-\frac2{n-1}W.
\]

(c) Apply Theorem~6.5 with $\mathcal G=\sigma(\pi)$ and $\lambda=\frac2{n-1}$. First,
$\mathbb E[D^2\mid\pi]=\frac1{n(n-1)}\sum_{i,j}(b_i-b_j)^2(c'_i-c'_j)^2$. Expanding both squares, the terms with a single factor $b_ib_j$ or $c'_ic'_j$ vanish because $\sum b=\sum c'=0$, and
\[
\sum_{i,j}(b_i-b_j)^2(c'_i-c'_j)^2=2n\sum_ib_i^2c'^2_i+2\sum_ib_i^2\sum_jc_j^2+4W^2.
\]
So $T=\frac{\mathbb E[D^2\mid\pi]}{2\lambda}=\frac12V+\frac{n-1}{2n}+\frac{W^2}n$ with $V=\sum_ib_i^2c_{\pi(i)}^2$, and $\mathbb ET=1$ because $\mathbb EV=\frac{n-1}n$. Thus
$\mathbb E|1-T|\le\frac12\sqrt{\operatorname{Var}V}+\frac{\mathbb E|W^2-1|}n\le\frac12\sqrt{\operatorname{Var}V}+\frac2n$. By (a) with $x_i=b_i^2$, $y_j=c_j^2$,
$\operatorname{Var}V\le\frac1{n-1}\sum b_i^4\sum c_j^4$ (centring decreases sums of squares). Second, $|b_i-b_j|^3\le4(|b_i|^3+|b_j|^3)$, so
$\sum_{i\ne j}|b_i-b_j|^3\le8n\sum|b|^3$, and since $I,J$ and $\pi$ are independent,
\[
\mathbb E|D|^3=\frac{\sum_{i\ne j}|b_i-b_j|^3}{n(n-1)}\cdot\frac{\sum_{k\ne l}|c_k-c_l|^3}{n(n-1)}\le\frac{64\sum|b|^3\sum|c|^3}{(n-1)^2},
\qquad\frac{\mathbb E|D|^3}{2\lambda}\le\frac{16}{n-1}\sum|b|^3\sum|c|^3.
\]
Theorem~6.5 gives $d_W\le2\mathbb E|1-T|+\frac{\mathbb E|D|^3}{2\lambda}$, which is the bound.

(d) Let $m=\frac{n+1}2$ and $Q=\sum_i(i-m)^2=\frac{n(n^2-1)}{12}$. Then $\rho=\frac1Q\sum_i(i-m)(\pi(i)-m)$, and with $b_i=c_i=(i-m)\,(n-1)^{1/4}Q^{-1/2}$ we get $\sum b^2\sum c^2=\sqrt{n-1}^2=n-1$ and
$W=\sqrt{n-1}\,\rho$. Now $|b_i|\le\frac n2(n-1)^{1/4}Q^{-1/2}=O(n^{-1/4})$, so $\sum b^4=O(n\cdot n^{-1})=O(1)$ and $\sum|b|^3=O(n\cdot n^{-3/4})=O(n^{1/4})$. All three
terms of the bound are $O(n^{-1/2})$; one finds numerically the bound equal to $28.8/\sqrt n$.

\tools{Theorem~6.5 (exchangeable pairs); Volume~4, Chapter~5.}
\leads Hoeffding proved the central limit theorem for these \emph{combinatorial} statistics $\sum_ia_{i\pi(i)}$ in 1951, and Bolthausen (1984) found the Berry--Esseen rate
by Stein's method. The constant in (c) is poor: the true distance for $n=20$ is already below $0.01$.
\end{solution}

\begin{solution}{pb:erdostaylor}
\idea Returns of the planar walk are rare, about $\frac1{\pi t}$ at time $t$, and the expected number up to $n$ grows like $\log n$. But the returns come in clusters:
once at home, the walk has the same chance as at the start to come back. The factorial moments of $N_n$ are those of an exponential variable, the law of a variable with
no memory.

\writeup (a) Rotate by $45^\circ$: $U_t=X_t+Y_t$ and $V_t=X_t-Y_t$ each move by $\pm1$ at every step, and the four steps of the planar walk correspond to the four pairs of
signs, so $U$ and $V$ are independent simple walks on $\Z$. Hence $P(S_{2t}=0)=P(U_{2t}=0)^2=\bigl(\binom{2t}t4^{-t}\bigr)^2=u(t)$. Since $\binom{2t}t4^{-t}\sim\frac1{\sqrt{\pi t}}$
(Volume~5, Corollary~3.6), $u(t)\sim\frac1{\pi t}$ and $G_n=\sum_{t\le n}u(t)\sim\frac1\pi\log n$.

(b) $N_n(N_n-1)\cdots(N_n-k+1)$ counts ordered $k$-tuples of distinct return times, so it equals $k!$ times the number of increasing ones. By the Markov property,
$P(S_{2t_1}=\dots=S_{2t_k}=0)=u(t_1)u(t_2-t_1)\cdots u(t_k-t_{k-1})$ for $t_1<\dots<t_k$. Hence
\[
\mathbb E[(N_n)_k]=k!\sum_{0<t_1<\dots<t_k\le n}\prod_{i}u(t_i-t_{i-1}).
\]
Writing $s_i=t_i-t_{i-1}\ge1$ (with $t_0=0$): dropping the constraint $\sum s_i\le n$ to $s_i\le n$ gives the upper bound $k!\,G_n^k$; requiring $s_i\le\lfloor n/k\rfloor$ for every $i$
implies $\sum s_i\le n$ and gives the lower bound $k!\,G_{\lfloor n/k\rfloor}^k$.

(c) Since $G_{\lfloor n/k\rfloor}/G_n\to1$, $\mathbb E[(N_n)_k]/G_n^k\to k!$ for every $k$. For a variable with values in $\N$ and $z\in[-1,0]$, the partial sums of
$(1+z)^N=\sum_k\binom Nkz^k$ alternately overestimate and underestimate it (Bonferroni's inequalities). With $z=-\frac s{G_n}$, $0<s<1$, and $n$ large:
\[
\sum_{k\le2K+1}\frac{(-s)^k}{k!}\frac{\mathbb E[(N_n)_k]}{G_n^k}\le\mathbb E\Bigl(1-\frac s{G_n}\Bigr)^{N_n}\le\sum_{k\le2K}\frac{(-s)^k}{k!}\frac{\mathbb E[(N_n)_k]}{G_n^k}.
\]
As $n\to\infty$ both sides tend to $\sum_{k\le2K+1}(-s)^k$ and $\sum_{k\le2K}(-s)^k$, then as $K\to\infty$ to $\frac1{1+s}$. Moreover $|(1-\frac s{G_n})^{N}-e^{-sN/G_n}|\to0$
uniformly in $N$, so $\mathbb Ee^{-sN_n/G_n}\to\frac1{1+s}$ for $0<s<1$. The sequence $N_n/G_n$ is tight (its mean is $1$); if $Y$ is the limit in law of a subsequence (Volume~7,
Theorem~11.7), then $\mathbb Ee^{-sY}=\frac1{1+s}$ for $0<s<1$. Both sides are holomorphic in $\operatorname{Re}s>0$, so they agree for all $s>0$ (Volume~5, Theorem~11.18). The
functions $y\mapsto e^{-ky}$, $k\in\N$, generate an algebra dense in the continuous functions on $[0,\infty]$ (Stone--Weierstrass), so the law of $Y$ is determined:
it is the exponential law, whose Laplace transform is $\frac1{1+s}$. Every subsequence has the same limit, so $N_n/G_n$ converges in law to it.

\tools{Volume~5, Corollary~3.6 (Wallis) and Theorem~11.18 (identity theorem); Volume~4, Chapter~7 (the Markov property, Pólya's theorem); Volume~7, Theorem~12.7 (Prokhorov).}
\leads Erdős and Taylor proved this in 1960. In dimension $1$ the number of returns up to $2n$, divided by $\sqrt n$, tends to a multiple of the absolute value of a Gaussian; in dimension
$3$ and more it stays bounded, with a geometric law. Dimension $2$ sits on the border, and the exponential law is the mark of a variable that keeps forgetting its past.
Numerically, the convergence is slow: for $n=10^4$, $\mathbb EN_n^2/(2G_n^2)=1.13$ instead of $1$.
\end{solution}

\begin{solution}{pb:dice}
\idea Modulo $q$, the total of $n$ dice is a random walk on $\Z/q\Z$, and its law is read off from the Fourier coefficients of one die. Modulo $7$ the six faces are all
the non-zero residues, which makes every coefficient equal to $-\frac16$.

\writeup (a) Let $\omega=e^{2\pi i/7}$ and $\phi(h)=\mathbb E\omega^{hX}=\frac16\sum_{f=1}^6\omega^{hf}$ for one die $X$. For $h\not\equiv0$, $\sum_{f=0}^6\omega^{hf}=0$, so
$\phi(h)=-\frac16$. By Fourier inversion on $\Z/7\Z$,
\[
P(S_n\equiv r)=\frac17\sum_{h=0}^6\omega^{-hr}\phi(h)^n=\frac17+\frac17\Bigl(-\frac16\Bigr)^n\sum_{h=1}^6\omega^{-hr}
=\begin{cases}\frac17+\frac67\bigl(-\frac16\bigr)^n,& r\equiv0,\\[2pt] \frac17-\frac17\bigl(-\frac16\bigr)^n,& r\not\equiv0.\end{cases}
\]
(b) Modulo $5$ the faces are $1,2,3,4,0,1$, so $\phi(h)=\frac16\bigl(\sum_{f=0}^4\omega^{hf}+\omega^h\bigr)=\frac{\omega^h}6$ for $h\not\equiv0$, with $\omega=e^{2\pi i/5}$. Then
$P(S_n\equiv r)=\frac15+\frac1{5\cdot6^n}\sum_{h=1}^4\omega^{h(n-r)}=\frac15+\frac{5\cdot\mathbf 1\{r\equiv n\}-1}{5\cdot6^n}$. Modulo $4$ the faces are $1,2,3,0,1,2$ and with
$\omega=i$: $\phi(1)=\frac{i-1}6$, $\phi(2)=\frac{-1+1-1+1-1+1}6=0$, $\phi(3)=\overline{\phi(1)}$. So
$P(S_n\equiv r)-\frac14=\frac12\operatorname{Re}\bigl(i^{-r}\phi(1)^n\bigr)$, and with $\theta_n=\frac{3\pi n}4$ the argument of $\phi(1)^n$,
$\|\mathcal L(S_n\bmod4)-U\|=\frac14\bigl(\frac{\sqrt2}6\bigr)^n\sum_{r=0}^3\bigl|\cos(\theta_n-\frac{\pi r}2)\bigr|=\frac12\bigl(\frac{\sqrt2}6\bigr)^n\bigl(|\cos\theta_n|+|\sin\theta_n|\bigr)$,
which is between $\frac12(\frac{\sqrt2}6)^n$ and $\frac1{\sqrt2}(\frac{\sqrt2}6)^n$. Modulo $2$, $\phi(1)=\frac16\sum_f(-1)^f=0$: three odd faces and three even ones, so a single die is
already uniform modulo $2$, and so is every total.

(c) The faces are $1+3J$ with $J$ uniform on $\{0,\dots,5\}$, so $\mu=1+3\cdot\frac52=\frac{17}2$ and $\sigma^2=9\cdot\frac{35}{12}=\frac{105}4$. The total $S_n=n+3\sum_iJ_i$ lives on
$n+3\Z$. The $J_i$ have span $1$, so the local limit theorem (Theorem~4.12) gives
$\sup_j\bigl|\frac\sigma3\sqrt n\,P(\sum J_i=j)-\phi\bigl(\frac{j-\frac52n}{\frac\sigma3\sqrt n}\bigr)\bigr|\to0$. With $k=n+3j$, $P(S_n=k)=P(\sum J_i=j)$ and
$\frac{j-\frac52n}{\sigma\sqrt n/3}=\frac{k-n\mu}{\sigma\sqrt n}$; multiplying by $3$ gives the claim. The factor $3$ is the span: the whole mass sits on one class modulo $3$,
so each possible value carries three times the Gaussian density.

\tools{Theorem~4.12 and Lemma~4.11 (local limit theorem, span); Volume~3, Chapter~11 (characters).}
\leads The rule ``the law modulo $q$ is uniform up to $\max_{h\ne0}|\phi(h)|^n$'' is the simplest case of the upper bound lemma used in the next problem.
\end{solution}

\begin{solution}{pb:squaresmodp}
\idea The upper bound lemma turns mixing into a statement about Fourier coefficients, and the Fourier coefficients of the law of a square are Gauss sums, all of
modulus $p^{-1/2}$. A $\pm1$ step has a coefficient $\cos\frac{2\pi}p$, very close to $1$: it mixes in $p^2$ steps, the squares in two.

\writeup (a) For $h\in\Z/p\Z$ let $\hat Q(h)=\mathbb Ee(ha^2/p)=\frac1p\sum_ae(ha^2/p)$, with $e(x)=e^{2\pi ix}$. For $h\ne0$, substituting $a=b+c$,
\[
\Bigl|\sum_ae\Bigl(\frac{ha^2}p\Bigr)\Bigr|^2=\sum_{a,b}e\Bigl(\frac{h(a^2-b^2)}p\Bigr)=\sum_c e\Bigl(\frac{hc^2}p\Bigr)\sum_be\Bigl(\frac{2hcb}p\Bigr)=p,
\]
because the inner sum is $p$ for $c=0$ and $0$ otherwise. So $|\hat Q(h)|=p^{-1/2}$, and the law $Q_k$ has coefficients $\hat Q(h)^k$. The upper bound lemma (Volume~10,
Lemma~10.4, for the abelian group $\Z/p\Z$) gives $\|Q_k-U\|^2\le\frac14\sum_{h\ne0}|\hat Q(h)|^{2k}=\frac{p-1}{4p^k}$.

(b) Let $N(r)$ be the number of $(a,b)$ with $a^2+b^2=r$. The number of $a$ with $a^2=x$ is $1+\left(\frac xp\right)$, so
$N(r)=\sum_{x+y=r}\bigl(1+(\frac xp)\bigr)\bigl(1+(\frac yp)\bigr)=p+\sum_x\left(\frac{x(r-x)}p\right)$, the linear terms summing to $0$. For $r=0$ the last sum is
$\sum_x\left(\frac{-x^2}p\right)=(p-1)\left(\frac{-1}p\right)$; for $r\ne0$, writing $x(r-x)=x^2(rx^{-1}-1)$ for $x\ne0$, it is $\sum_{y\ne-1}\left(\frac yp\right)=-\left(\frac{-1}p\right)$.
With $\eta=\left(\frac{-1}p\right)$: $Q_2(0)-\frac1p=\frac{(p-1)\eta}{p^2}$ and $Q_2(r)-\frac1p=-\frac\eta{p^2}$ for $r\ne0$, so
$\|Q_2-U\|=\frac12\bigl(\frac{p-1}{p^2}+\frac{p-1}{p^2}\bigr)=\frac{p-1}{p^2}$. The bound of (a) is $\frac{\sqrt{p-1}}{2p}$, about $\frac1{2\sqrt p}$ instead of $\frac1p$: right in spirit, off
by $\sqrt p$.

(c) For any law $Q$ and $h\ne0$, $|\hat Q(h)|=\bigl|\sum_x(Q(x)-\frac1p)e(hx/p)\bigr|\le2\|Q-U\|$. For the $\pm1$ walk, $\hat Q_k(1)=\cos^k\frac{2\pi}p$, so
$\|Q_k-U\|\ge\frac12\cos^k\frac{2\pi}p$. Since $\log\cos u\ge-\frac{u^2}2-\frac{u^4}6$ for $|u|\le1$, for $k\le p^2/10$,
\[
\|Q_k-U\|\ge\frac12\exp\Bigl(-\frac{p^2}{10}\Bigl(\frac{2\pi^2}{p^2}+\frac{8\pi^4}{3p^4}\Bigr)\Bigr)\longrightarrow\frac12e^{-\pi^2/5}=0.069>\frac1{20}.
\]
The truth is worse: one finds numerically $\|Q_k-U\|=0.40$ for $k=p^2/10$, because the frequency $h=\frac{p-1}2$, with $\cos\frac{2\pi h}p$ close to $-1$, also sees that the
walk remembers its parity.
Squares mix fast because the law of $a^2$ is spread evenly over all frequencies: no character sees it as concentrated. The $\pm1$ steps are concentrated near $0$, and
the low frequencies see that for a long time.

\tools{Lemma~14.4 (upper bound lemma); Volume~2, Chapter~6 (Legendre symbol); Volume~2, Exercise~6.7 (Gauss sums).}
\leads The same computation with $a^d$ instead of $a^2$ needs the bound $|\sum_ae(ha^d/p)|\le(d-1)\sqrt p$, which follows from Gauss sums of multiplicative
characters: then $d$-th powers mix in a bounded
number of steps when $p$ is large, a probabilistic shadow of Waring's problem modulo $p$.
\end{solution}

\begin{solution}{pb:lucky}
\idea Counting lucky tickets is computing $P(D_1+\dots+D_m=0)$ for the differences $D_i$ of two digits, a symmetric walk on $\Z$. The local limit theorem gives the main
term, and one more term of the expansion of $\log\mathbb Ee^{itD}$ gives the correction.

\writeup (a) Let $N_m(s)$ be the number of strings of $m$ digits with sum $s$, so that $L_m=\sum_sN_m(s)^2$ and $\sum_sN_m(s)z^s=P(z)^m$ with $P(z)=1+z+\dots+z^9$. By
Parseval's formula for trigonometric polynomials, $L_m=\frac1{2\pi}\int_{-\pi}^\pi|P(e^{it})|^{2m}\dd t$, and $|P(e^{it})|=\bigl|\frac{e^{10it}-1}{e^{it}-1}\bigr|=\bigl|\frac{\sin5t}{\sin\frac t2}\bigr|$.

(b) Equivalently $L_m=10^{2m}P(D_1+\dots+D_m=0)$ where $D_i=X_i-Y_i$ is the difference of two independent uniform digits: symmetric, with values in $\{-9,\dots,9\}$,
span $1$, and variance $\sigma^2=2\cdot\frac{99}{12}=\frac{33}2$. The local limit theorem (Theorem~4.12) gives
$P(\sum D_i=0)\sim\frac1{\sqrt{2\pi m\sigma^2}}=\frac1{\sqrt{33\pi m}}$.

(c) Let $\psi(t)=\mathbb Ee^{itD}=\frac{|P(e^{it})|^2}{100}$, real and even, with $|\psi(t)|<1$ for $0<|t|\le\pi$, so that $|\psi|\le1-c$ on $\delta\le|t|\le\pi$: that part of
$\frac1{2\pi}\int\psi^m$ is exponentially small. Near $0$, $\log\psi(t)=-\frac{\sigma^2t^2}2+\frac{\kappa_4t^4}{24}+O(t^6)$, where $\kappa_4=2\kappa_4(X)$ is the fourth cumulant
of $D$ ($X$ and $-Y$ contribute equally). For $X$ uniform on $\{0,\dots,9\}$, $\sigma_X^2=\frac{33}4$ and $\mathbb E(X-\frac92)^4=\frac15\sum_{j=0}^4(j+\frac12)^4=\frac{9669}{80}$, so
$\kappa_4(X)=\frac{9669}{80}-3\cdot\frac{1089}{16}=-\frac{3333}{40}$ and $\kappa_4=-\frac{3333}{20}=-\frac{9999}{60}$. Substituting
$t=u/\sqrt m$,
\[
\frac1{2\pi}\int_{|t|<\delta}\psi^m\dd t=\frac1{2\pi\sqrt m}\int e^{-\sigma^2u^2/2}\Bigl(1+\frac{\kappa_4u^4}{24m}+O\Bigl(\frac{u^6+u^8}{m^2}\Bigr)\Bigr)\dd u+(\text{exponentially small}),
\]
and $\int u^4e^{-\sigma^2u^2/2}\dd u=\frac{3\sqrt{2\pi}}{\sigma^5}$. So $P(\sum D_i=0)=\frac1{\sigma\sqrt{2\pi m}}\bigl(1+\frac{\kappa_4}{8\sigma^4m}+O(m^{-2})\bigr)$, and
$\frac{\kappa_4}{8\sigma^4}=-\frac{9999/60}{8\cdot1089/4}=-\frac{9999}{130\,680}=-\frac{101}{1320}$. For $m=3$: $\frac{10^6}{\sqrt{99\pi}}=56\,703$ and the corrected value $55\,257$,
against $55\,252$ lucky tickets.

\tools{Theorem~4.12 (local limit theorem); Volume~5, Theorem~8.11 (Parseval); Volume~4, Chapter~6 (characteristic functions).}
\leads The expansion continues into an asymptotic series in powers of $\frac1m$, the Edgeworth expansion of the local limit theorem; the same
computation counts the strings of $2m$ digits in any base whose two halves have equal sums.
\end{solution}

\begin{solution}{pb:signedsums}
\idea The same Fourier integral as in the local limit theorem, but for non-identical summands $\pm k$. The integrand $\prod\cos(kt)$ has modulus $1$ at $t=0$ and
also at $t=\pm\pi$, and the second peak doubles the answer.

\writeup (a) The sum $\pm1\pm\dots\pm n$ has the parity of $\frac{n(n+1)}2$, which is even if and only if $n\equiv0,3\pmod4$. And
$P(\sum_k\varepsilon_kk=0)=\frac1{2\pi}\int_{-\pi}^\pi\mathbb Ee^{it\sum\varepsilon_kk}\dd t=\frac1{2\pi}\int_{-\pi}^\pi\prod_k\cos(kt)\dd t$.

(b) Let $s_n^2=\sum_kk^2=\frac{n(n+1)(2n+1)}6\sim\frac{n^3}3$ and $g(t)=\prod_{k\le n}\cos(kt)$. Since $|\cos x|\le e^{-\sin^2x/2}$ and $\sum_{k\le n}\sin^2(kt)=\frac n2-\frac12\sum_{k\le n}\cos(2kt)\ge\frac n2-\frac1{2|\sin t|}$,
we get $|g(t)|\le e^{-3n/16}$ whenever $|\sin t|\ge\frac4n$. The remaining $t$ are near $0$ or $\pm\pi$, and $g(t\pm\pi)=(-1)^{n(n+1)/2}g(t)=g(t)$, so it suffices to study $|t|\le\frac7n$ (if $|t|\le\frac\pi2$ and $|\sin t|<\frac4n$, then $|t|<\frac{2\pi}n$). For $\frac1{n}\le|t|\le\frac7n$, the factors with $k\le\frac1{|t|}$ satisfy $|kt|\le1$ and $\cos(kt)\le e^{-k^2t^2/2}$, so
$|g(t)|\le e^{-t^2\sum_{k\le1/|t|}k^2/2}\le e^{-1/(48|t|)}\le e^{-n/336}$. For $|t|\le\frac1n$, all $|kt|\le1$ and $\log\cos x=-\frac{x^2}2+O(x^4)$, so
$g(t)=e^{-s_n^2t^2/2+O(t^4n^5)}$ and $|g(t)|\le e^{-s_n^2t^2/2}$. Split at $|t|=n^{-3/2}\log n$: inside, $t^4n^5=O(n^{-1}\log^4n)$ and $\int e^{-s_n^2t^2/2}\dd t=\frac{\sqrt{2\pi}}{s_n}(1+o(1))$;
outside, $|g|\le e^{-c\log^2n}$. Hence $\int_{-\pi}^\pi g=2\frac{\sqrt{2\pi}}{s_n}(1+o(1))$, the factor $2$ coming from the two peaks, and
\[
Z(n)=\frac{2^n}{2\pi}\cdot\frac{2\sqrt{2\pi}}{s_n}(1+o(1))\sim2^n\sqrt{\frac2\pi}\cdot\frac{\sqrt3}{n^{3/2}}=\sqrt{\frac6\pi}\,\frac{2^n}{n^{3/2}}.
\]

(c) A choice of signs with sum $0$ is the same as the set $S$ of plus signs, with $\sum_{k\in S}k=\frac{n(n+1)}4$: there are $Z(n)$ such subsets, and $\frac{Z(n)}2$ unordered
splittings of $\{1,\dots,n\}$ into two parts of equal sum.

\tools{Volume~11, Chapter~4 (the Fourier proof of the local limit theorem); Volume~5, Chapter~8 (Dirichlet kernel).}
\leads The answer to the opening question is the middle one: $2^n/n^{3/2}$. The sequence $Z(n)$ is $0,0,2,2,0,0,8,14,0,0,70,124,\dots$; van Lint (1967) and Entringer (1968)
studied it, and the same method counts the subsets of $\{1,\dots,n\}$ with a given sum modulo $q$.
\end{solution}

\begin{solution}{pb:smirnov}
\idea The event ``the path stays in the strip'' is the complement of a union of events ``the path visits $h$, then $-h$, then $h$\dots'', and each of those is counted by
reflecting the path several times. The result is an alternating sum of binomial coefficients, which converges term by term to an alternating sum of Gaussians.

\writeup (a) For a path $S$ from $0$ to $0$ of length $2n$ that reaches $\{h,-h\}$, write down the sequence of barriers it visits, keeping only the changes (visits to $h$
followed by a visit to $-h$ followed by $h$\dots): this \emph{reduced sequence} alternates, starts with $h$ or $-h$, and has some length $\ell\ge1$. Let $A_j$ be the set of
paths whose reduced sequence contains $h,-h,h,\dots$ ($j$ terms) as an initial alternation, that is, which visit $h$ at some time $t_1$, then $-h$ at some $t_2>t_1$, and
so on $j$ times; define $B_j$ in the same way starting with $-h$. A path with reduced sequence of length $\ell$ starting with $h$ lies in $A_j$ exactly for $j\le\ell$
and in $B_j$ exactly for $j\le\ell-1$, so $\sum_{j\ge1}(-1)^{j-1}\bigl(\mathbf 1_{A_j}+\mathbf 1_{B_j}\bigr)=\mathbf 1\{\ell\text{ odd}\}+\mathbf 1\{\ell-1\text{ odd}\}=1$. Hence the number of
paths that reach the barriers is $\sum_{j\ge1}(-1)^{j-1}(|A_j|+|B_j|)$.

To count $A_j$, reflect the part of the path after its first visit to $h$ about $h$ (Figure~\ref{fig:strip}): the end point $0$ goes to $2h$, and a later visit to $-h$
becomes a visit to $3h$. Reflect again after the first visit to $3h$ about $3h$, and so on. After $j$ reflections the path ends at $2jh$, and the construction is a bijection
between $A_j$ and all paths from $0$ to $2jh$ (any such path reaches $h,3h,\dots,(2j-1)h$ in this order, and the reflections can be undone). So
$|A_j|=\binom{2n}{n+jh}$, and $|B_j|$ is the same by symmetry. The number of paths in the strip is
\[
\binom{2n}n-2\sum_{j\ge1}(-1)^{j-1}\binom{2n}{n+jh}=\sum_{k\in\Z}(-1)^k\binom{2n}{n+kh}.
\]

\begin{figure}[ht]
\centering\fig{ch01-strip}
\caption{A path of length $40$ from $0$ to $0$ that visits $h=4$ (charcoal). Reflecting it about $h$ after its first visit (dotted rust) gives a path from $0$ to $2h$:
the first step of the inclusion--exclusion count of Problem~\ref{pb:smirnov}.}
\label{fig:strip}
\end{figure}

(b) $\binom{2n}{n+k}\big/\binom{2n}n=\prod_{i=1}^k\frac{n-i+1}{n+i}$ for $k\ge0$. Each factor is $1-\frac{2i-1}{n+i}\le e^{-(2i-1)/(2n)}$, so the ratio is at most $e^{-k^2/(2n)}$. And
$\log\frac{n-i+1}{n+i}=\log\bigl(1-\frac{i-1}n\bigr)-\log\bigl(1+\frac in\bigr)=-\frac{2i-1}n+\frac{2i-1}{2n^2}+O\bigl(\frac{i^3}{n^3}\bigr)$, whose sum over $i\le k$ is
$-\frac{k^2}n+O\bigl(\frac{k^2}{n^2}+\frac{k^4}{n^3}\bigr)$.

(c) $D_n=\frac1n\max_k|S_k|$, so $\sqrt{n/2}\,D_n<x$ means $\max_k|S_k|<x\sqrt{2n}$, that is, the path stays strictly inside the strip of half-width $h_n=\lceil x\sqrt{2n}\rceil$.
By (a), $P(\sqrt{n/2}\,D_n<x)=\sum_k(-1)^k\binom{2n}{n+kh_n}\big/\binom{2n}n$. By (b), the $k$-th term tends to $(-1)^ke^{-2k^2x^2}$, since $h_n^2/n\to2x^2$, and it is bounded by
$e^{-k^2h_n^2/(2n)}\le e^{-k^2x^2}$, which is summable. By dominated convergence the sum tends to $\sum_k(-1)^ke^{-2k^2x^2}=K(x)$.

(d) Apply Poisson summation (Volume~5, Theorem~10.9) to $f(u)=e^{-2x^2u^2+i\pi u}$, whose values at integers are $(-1)^ke^{-2k^2x^2}$ and whose Fourier transform is
$\hat f(m)=\int f(u)e^{-2\pi imu}\dd u=\sqrt{\frac\pi{2x^2}}\,e^{-\pi^2(2m-1)^2/(8x^2)}$. Each odd number $2m-1$, positive or negative, occurs once, so
$K(x)=\sum_m\hat f(m)=\frac{\sqrt{2\pi}}x\sum_{k\ge1}e^{-(2k-1)^2\pi^2/(8x^2)}$. At $x=1.36$: $1-2e^{-2x^2}=1-2(0.02474)=0.9505$, with a next term of $3.7\times10^{-7}$; on the
right, the first term is $0.9459$ and the second $0.0046$. For $x=0.3$ the left series starts with $1-2e^{-0.18}=-0.67$ and converges slowly, while one term on the right gives
$K(0.3)=9.3\times10^{-6}$: the left series is for large $x$, the right one for small $x$.

(e) Numerically, one finds $\sqrt N\sup_t|F_N(t)-t|=0.0265$ for $N=10^4$. The usual test rejects for \emph{large} values ($1.36$ at level $5\%$), so it accepts. But the probability
that $N$ random points do as well is about $K(0.0265)$, which by the right-hand series is $10^{-761}$: the sequence is far too regular to be random. A good test looks at both
tails, and an irrational rotation fails it spectacularly.

\tools{Volume~4, Lemma~1.11 (reflection principle); Volume~5, Theorem~10.9 (Poisson summation); Volume~11, Chapter~3 (Kolmogorov's statistic).}
\leads Smirnov proved (c) in 1939, Kolmogorov the one-sample version in 1933; Gnedenko and Korolyuk (1951) found the reflection proof above. The identity of (d) is a
case of Jacobi's transformation of theta functions, and it says that the maximum of a Brownian bridge has the same law whether one looks at it through its
paths or through heat diffusion in a strip.
\end{solution}

\begin{solution}{pb:primeseries}
\idea For independent bounded symmetric terms, the three series reduce to one: the sum of the variances. Over the primes, Mertens's theorem says that $\sum\frac1p$ grows like
$\log\log$, so a single power of $\log p$ is already too much.

\writeup (a) The terms $X_p=\frac{\varepsilon_p}{\sqrt p(\log p)^a}$ are independent, symmetric and bounded (they tend to $0$). With $c$ larger than all of them, the three
series of Theorem~2.5 are $0$, $0$ and $\sum_p\frac1{p(\log p)^{2a}}$, so the series converges almost surely if and only if this sum is finite, and
diverges almost surely otherwise. Group the primes in blocks $e^{2^j}<p\le e^{2^{j+1}}$. By Mertens's theorem (Volume~6, Theorem~1.5),
$\sum_{\text{block}}\frac1p=\log2^{j+1}-\log2^j+o(1)\to\log2$, and $(\log p)^{2a}$ is between $2^{2aj}$ and $2^{2a(j+1)}$ on the block. So the sum converges if and only if
$\sum_j2^{-2aj}<\infty$, that is, $a>0$.

(b) The same theorem over the integers: $\sum_n\frac1{n(\log n)^{2a}}<\infty$ if and only if $2a>1$, by comparison with $\int\frac{\dd t}{t(\log t)^{2a}}$.

(c) Let $S_N$ be the partial sums and $s_N^2=\operatorname{Var}S_N\to\infty$. The summands are bounded, so Lindeberg's condition holds and $S_N/s_N\Rightarrow N(0,1)$; for fixed $M$,
$P(S_N>M)\to\frac12$. Hence $P(S_N>M\text{ for infinitely many }N)\ge\limsup_NP(S_N>M)=\frac12$. The event $\{\limsup S_N=+\infty\}$ is a tail event (changing finitely many terms
changes all $S_N$ by a bounded amount), it contains the decreasing limit of the events above as $M\to\infty$, so it has probability at least $\frac12$, hence $1$ by
Kolmogorov's $0$--$1$ law (Theorem~1.4). By symmetry $\liminf S_N=-\infty$ too.

(d) $\sum_pp^{-2\sigma}<\infty$ for $\sigma>\frac12$, and at $\sigma=\frac12$ we are in case $a=0$ of (a).

\tools{Theorems~1.4 and~2.5 (zero-one law, three series); Volume~6, Theorem~1.5 (Mertens); Volume~4, Theorem~6.10.}
\leads The random Euler product $\prod_p(1-\varepsilon_pp^{-s})^{-1}$ is a model of $\zeta(s)$ near the critical line, used by Bohr and Jessen (1930) to describe the values of
$\log\zeta(\sigma+it)$ for fixed $\sigma>\frac12$. At $\sigma=\frac12$ the random model diverges, as (c) shows, and so does the analogy.
\end{solution}

\begin{solution}{pb:factorialseries}
\idea The terms decrease so fast that each one dominates the sum of all the later ones: the first sign that differs decides which sum is larger. So $X$ codes its signs
injectively, as binary digits do, but with gaps that shrink superexponentially.

\writeup Let $r_n=\sum_{k>n}\frac1{k!}$. Then $r_n=\frac1{(n+1)!}\bigl(1+\frac1{n+2}+\frac1{(n+2)(n+3)}+\dots\bigr)<\frac1{(n+1)!}\cdot\frac{n+1}n=\frac1{n\cdot n!}$, which is at most $\frac1{n!}$.

(a) If two sign sequences first differ at $n$, the two sums differ by $\frac2{n!}$ plus a quantity of absolute value at most $2r_n<\frac2{n!}$: they are different.

(b) By (a), $X=x$ for at most one sign sequence, and every single sequence has probability $\lim_n2^{-n}=0$. So $P(X=x)=0$ for all $x$.

(c) $X$ is a continuous function of the sequence of signs, so its range $C$ is compact (the image of $\{\pm1\}^{\N}$). Given the first $n$ signs, with partial sum $s$,
$X\in[s-r_n,s+r_n]$: the $2^n$ choices of the first $n$ signs give $2^n$ intervals of length $2r_n\le\frac2{n\cdot n!}$ covering $C$.

(d) For $s>0$, $\mathcal H^s(C)\le\liminf_n2^n\bigl(\frac2{n\cdot n!}\bigr)^s=0$, because $n!$ grows faster than any $c^n$. So $\dim_HC=0$ (Volume~7, Chapter~12). In particular $C$ is
Lebesgue-null: the law of $X$ has no atom by (b) and is carried by a null set, so it is continuous and singular.

\tools{Volume~7, Chapter~12 (Hausdorff measures and dimension) and Corollary~11.6.}
\leads Compare $\sum\varepsilon_n2^{-n}$, uniform, and $\sum\varepsilon_n3^{-n}$, carried by a Cantor set of dimension $\frac{\log2}{\log3}$: the dimension of the law of
$\sum\varepsilon_na_n$ is governed by how fast $a_n$ decreases, and superexponential decay gives dimension $0$.
\end{solution}

\begin{solution}{pb:smooth}
\idea A density is smooth when its characteristic function decays faster than any power. A product of cosines decays when many factors are visibly less than $1$:
for $\frac t{n^2}$, the factors with $n\approx\sqrt t$ are the useful ones, and there are enough of them to give $e^{-c\sqrt t}$.

\writeup (a) The partial sums $X_N$ converge almost surely (three series, $\sum n^{-4}<\infty$), hence in law, so $\varphi(t)=\lim_N\mathbb Ee^{itX_N}=\lim_N\prod_{n\le N}\cos\frac t{n^2}$
by independence (Theorem~3.3).

(b) $\log\cos u=-\frac{u^2}2-\frac{u^4}{12}-\dots\le-\frac{u^2}2$ for $|u|<\frac\pi2$ (all the coefficients of the series are negative). Let $|t|\ge1$ and $N=\lceil\sqrt{|t|}\rceil\le2\sqrt{|t|}$. For
$n\ge N$, $\frac{|t|}{n^2}\le1$, and the other factors have modulus at most $1$, so
\[
|\varphi(t)|\le\exp\Bigl(-\frac{t^2}2\sum_{n\ge N}\frac1{n^4}\Bigr)\le\exp\Bigl(-\frac{t^2}2\cdot\frac1{3N^3}\Bigr)\le\exp\Bigl(-\frac{t^2}{48|t|^{3/2}}\Bigr)=e^{-\sqrt{|t|}/48},
\]
using $\sum_{n\ge N}n^{-4}\ge\int_N^\infty x^{-4}\dd x$.

(c) For every $k$, $|t|^k|\varphi(t)|$ is integrable. By the inversion theorem (Volume~8, Theorem~18.2, or Volume~4, Theorem~6.8), $X$ has the continuous density
$f(x)=\frac1{2\pi}\int e^{-itx}\varphi(t)\dd t$, and differentiation under the integral sign, dominated by $|t|^k|\varphi(t)|$, shows that $f$ is $C^k$ for every $k$.

(d) The series converges almost surely because $\sum n^{-2\alpha}<\infty$. With $N=\lceil|t|^{1/\alpha}\rceil\le2|t|^{1/\alpha}$, the same computation gives
$\sum_{n\ge N}n^{-2\alpha}\ge\frac{N^{1-2\alpha}}{2\alpha-1}\ge\frac{2^{1-2\alpha}|t|^{1/\alpha-2}}{2\alpha-1}$ and
$|\varphi(t)|\le\exp\bigl(-\frac{2^{-2\alpha}}{2\alpha-1}|t|^{1/\alpha}\bigr)$, so the density is again $C^\infty$. For $\sum\varepsilon_n2^{-n}$ the characteristic function is
$\frac{\sin t}t$, which decays only like $\frac1{|t|}$: the uniform density is not even continuous. Geometric decay of the pieces is too fast for smoothness, polynomial decay is
slow enough.

\tools{Theorems~2.5 and~3.3; Volume~8, Theorem~18.2 (Fourier inversion).}
\leads Direct computation gives $-\log|\varphi(t)|/\sqrt t$: it is $1.06$ at $t=10^3$ and $1.01$ at $t=10^4$, so the true decay is about $e^{-\sqrt t}$ and the constant $\frac1{48}$ is
generous. For $\alpha\le\frac12$ the series diverges almost surely (three series again).
\end{solution}

\begin{solution}{pb:purity}
\idea Changing finitely many of the $X_n$ moves $X$ by an element of a countable group $H$. So for any set $E$, ``$X\in E+H$'' is a tail event, of probability $0$ or $1$.
Taking for $E$ an atom, or a null set charged by $\mu$, forces $\mu$ to choose one type.

\writeup (a) Let $V_n$ be the countable set of values of $X_n$ and $H$ the additive group generated by $\bigcup_nV_n$, which is countable. Almost surely every $X_n$ lies
in $V_n$; on that event, $\sum_{n\le m}X_n\in H$ for every $m$, so for a Borel set $E$ the set $E+H=\bigcup_{h\in H}(E+h)$ is Borel and
$\{X\in E+H\}=\bigl\{\sum_{n>m}X_n\in E+H\bigr\}$ for every $m$. So $\{X\in E+H\}$ coincides almost surely with a tail event, and $\mu(E+H)\in\{0,1\}$ by Kolmogorov's $0$--$1$
law (Theorem~1.4).

If $\mu$ has an atom $a$, then $\mu(a+H)\ge\mu(\{a\})>0$, so $\mu(a+H)=1$: $\mu$ is carried by the countable set $a+H$. Otherwise $\mu$ is continuous. If it is not absolutely
continuous, there is a Lebesgue-null Borel set $E$ with $\mu(E)>0$; then $E+H$ is null (a countable union of translates of $E$) and $\mu(E+H)=1$, so $\mu$ is carried by a
null set: it is singular. By the Lebesgue decomposition into three parts (Volume~7, Corollary~11.6), these are the three pure types, and we have shown that only one part can
be non-zero.

(b) The atoms of $\mu_\lambda$, the law of $X_\lambda$, have total mass at most $1$, so if there is one, the largest mass $M>0$ is attained, by a finite non-empty set $A$
of atoms. Write $X_\lambda=\lambda\varepsilon_1+\lambda X'$, where $X'=\sum_{n\ge1}\varepsilon_{n+1}\lambda^n$ has the same law as $X_\lambda$ and is independent of $\varepsilon_1$. For
$a\in A$,
\[
M=P(X_\lambda=a)=\tfrac12P\bigl(X'=\tfrac a\lambda-1\bigr)+\tfrac12P\bigl(X'=\tfrac a\lambda+1\bigr),
\]
and both probabilities are at most $M$, so both equal $M$: $\frac a\lambda\pm1\in A$. Apply this to $a=\max A$: $\frac a\lambda+1\le a$, so $a(\frac1\lambda-1)\le-1$ and $a<0$. Apply it to
$b=\min A$: $\frac b\lambda-1\ge b$, so $b>0$. But $b\le a$: contradiction. So $\mu_\lambda$ has no atom.

(c) By (a) and (b), $\mu_\lambda$ is either absolutely continuous or singular. For $\lambda=\frac12$, $X_\lambda=\sum\varepsilon_n2^{-n}$ is uniform on $[-1,1]$ (binary digits),
absolutely continuous. For $\lambda=\frac13$, $X_\lambda$ lies, given its first $n$ signs, in an interval of length $2\sum_{k>n}3^{-k}=3^{-n}$: the $2^n$ intervals have total length
$(\frac23)^n\to0$, so $\mu_{1/3}$ is carried by a null set (a Cantor set) and is singular. The same argument works for all $\lambda<\frac12$.

\tools{Theorem~1.4 (zero-one law); Volume~7, Theorem~11.5 and Corollary~11.6 (Lebesgue decomposition).}
\leads This is the Jessen--Wintner purity theorem (1935). Lévy had shown that such a law is continuous if and only if $\prod_n\max_xP(X_n=x)=0$. For $\frac12<\lambda<1$
the type of $\mu_\lambda$ is the Bernoulli convolution problem of the next problem.
\end{solution}

\begin{solution}{pb:golden}
\idea If $\mu$ had a density, its Fourier transform would tend to $0$ (Riemann--Lebesgue). At the frequencies $\pi\varphi^k$ it does not: the cosines $\cos(\pi\varphi^m)$
are all close to $\pm1$, because $\varphi^m$ is extraordinarily close to an integer, the Lucas number $L_m$.

\writeup (a) $\varphi$ and $\psi=-\frac1\varphi$ are the roots of $x^2=x+1$, so $L_m=\varphi^m+\psi^m$ satisfies $L_{m+2}=L_{m+1}+L_m$, with $L_0=2$, $L_1=1$: all integers.

(b) As in Problem~\ref{pb:smooth}, $\hat\mu(t)=\prod_{n\ge1}\cos(t\varphi^{-n})$. At $t=\pi\varphi^k$ the exponents $k-n$ run over $k-1,k-2,\dots,0,-1,-2,\dots$:
\[
\hat\mu(\pi\varphi^k)=\prod_{m=0}^{k-1}\cos(\pi\varphi^m)\cdot\prod_{j\ge1}\cos(\pi\varphi^{-j}).
\]
Since $\pi\varphi^m=\pi L_m-\pi\psi^m$, $|\cos(\pi\varphi^m)|=|\cos(\pi\psi^m)|=|\cos(\pi\varphi^{-m})|$. So
$|\hat\mu(\pi\varphi^k)|=\prod_{m=1}^{k-1}|\cos(\pi\varphi^{-m})|\cdot c\to c^2$ (the factor $m=0$ is $|\cos\pi|=1$). Finally $c>0$: no factor vanishes, because $\varphi^{-j}$ is
irrational and so never in $\frac12+\Z$, and $\sum_j\bigl(1-|\cos(\pi\varphi^{-j})|\bigr)\le\sum_j\frac{\pi^2}2\varphi^{-2j}<\infty$, so the infinite product converges to a non-zero limit.
Numerically $c=0.0813$ and $c^2=0.00661$.

(c) If $\mu$ had a density $g\in L^1$, then $\hat\mu(t)=\int e^{itx}g(x)\dd x\to0$ as $|t|\to\infty$ by the Riemann--Lebesgue lemma (Volume~8, Proposition~18.1(b)); by (b) it does not.
By Problem~\ref{pb:purity}(b) and (a), $\mu$ has no atom and is pure, so it is continuous and singular. All the $2^n$ partial sums lie in $[-\varphi,\varphi]$, since
$\sum_{n\ge1}\varphi^{-n}=\frac1{\varphi-1}=\varphi$, and they overlap heavily (Figure~\ref{fig:golden}): the law looks smooth, and it is not.

\begin{figure}[ht]
\centering\fig{ch01-golden}
\caption{Histograms of $2\times10^6$ samples of $\sum_{n\le44}\varepsilon_n\lambda^n$, with $200$ bins. For $\lambda=\frac1\varphi$ (top) the law is singular (Problem~\ref{pb:golden}),
and the histogram is visibly rough; for $\lambda=0.7$ (bottom) Solomyak's theorem says that a density exists for almost every $\lambda$, and the picture agrees.}
\label{fig:golden}
\end{figure}

(d) Let $\theta_2,\dots,\theta_d$ be the other conjugates of $\theta$, $|\theta_i|<1$, and $\operatorname{tr}(\theta^m)=\theta^m+\sum_i\theta_i^m\in\Z$ (a symmetric function of the
roots of a monic integer polynomial). Then $|\cos(\pi\theta^m)|=|\cos(\pi\sum_i\theta_i^m)|$, and $\bigl|\sum_i\theta_i^m\bigr|\le(d-1)r^m$ with $r=\max|\theta_i|<1$, so
$\sum_m\bigl(1-|\cos(\pi\theta^m)|\bigr)<\infty$. No factor vanishes: $\theta^m\in\frac12+\Z$ is impossible for an algebraic integer, and $\theta^{-j}\in\frac12+\Z$ would make $\theta^j$
rational, hence an integer $N$; then each conjugate $\theta_i$ satisfies $\theta_i^j=N$ too and $|\theta_i|=N^{1/j}=\theta>1$, a contradiction as $d\ge2$. So
$|\hat\mu(\pi\theta^k)|=\prod_{m=0}^{k-1}|\cos(\pi\theta^m)|\cdot\prod_{j\ge1}|\cos(\pi\theta^{-j})|$ tends to a non-zero limit, and the law is singular as in (c).

\tools{Volume~8, Proposition~18.1(b) (Riemann--Lebesgue); Problem~\ref{pb:purity}; Volume~13, Chapter~3 (algebraic integers).}
\leads Erdős found this in 1939. Salem (1944) proved the converse for the Fourier side: $\hat\mu_\lambda(t)\not\to0$ only when $\frac1\lambda$ is a Pisot number. Whether
$\mu_\lambda$ has a density for every other $\lambda\in\left]\frac12,1\right[$ is still open: Solomyak (1995) proved it for almost every $\lambda$, Shmerkin (2014) showed that
the exceptional set has Hausdorff dimension $0$, and Varjú (2019) proved it for many algebraic $\lambda$ close to $1$.
\end{solution}

\begin{solution}{pb:kaczeros}
\idea Count the zeros of $f$ by the time it spends near $0$, weighted by its speed. Taking expectations turns the count into an integral of the density of $f(t)$ at $0$
times the mean speed given $f(t)=0$: Gaussian computations. The answer has a beautiful closed form, and almost all the real roots live near $\pm1$.

\writeup (a) The zeros of $g$ are simple, hence isolated, hence finitely many in $[a,b]$; call them $x_1<\dots<x_N$, all interior. Around each $x_i$ choose an interval
$I_i$ where $g'\ne0$, so that $g$ is strictly monotone on $I_i$ and $g$ has no zero outside $\bigcup I_i$. Let $\eta>0$ be the minimum of $|g|$ on $[a,b]\setminus\bigcup I_i$ and
of $|g|$ at the endpoints of the $I_i$. For $\varepsilon<\eta$, $\{|g|<\varepsilon\}\subseteq\bigcup I_i$ and on each $I_i$ the change of variables $y=g(t)$ gives
$\int_{I_i}\mathbf 1_{|g|<\varepsilon}|g'|=2\varepsilon$. So the integral is $2\varepsilon N$.

(b) $f$ has a multiple root if and only if the discriminant of $f$ vanishes (or $a_n=0$), a non-zero polynomial equation in $(a_0,\dots,a_n)$; its zero set is Lebesgue-null, so it
has probability $0$. Almost surely $f'$ is a non-zero polynomial of degree at most $n-1$, so $[a,b]$ splits into at most $n$ intervals on which $f$ is monotone, and on each,
$\int\mathbf 1_{|f|<\varepsilon}|f'|$ is the length of $f(I)\cap\left]-\varepsilon,\varepsilon\right[$, at most $2\varepsilon$.

(c) Almost surely $f(a)f(b)\ne0$, so by (a) the number $N[a,b]$ of roots in $[a,b]$ is the limit of $\frac1{2\varepsilon}\int_a^b\mathbf 1_{|f|<\varepsilon}|f'|$, which is bounded by $n$ by
(b). By dominated convergence and Fubini,
$\mathbb EN[a,b]=\lim_{\varepsilon\to0}\int_a^b\frac1{2\varepsilon}\mathbb E\bigl[\mathbf 1_{|f(t)|<\varepsilon}|f'(t)|\bigr]\dd t$. Fix $t$. $(f(t),f'(t))$ is a centred Gaussian vector with
$\operatorname{Var}f(t)=A$, $\operatorname{Cov}(f(t),f'(t))=B$, $\operatorname{Var}f'(t)=C$; given $f(t)=y$, $f'(t)$ is Gaussian with mean $\frac{By}A$ and variance $C-\frac{B^2}A$. So
\[
\begin{aligned}
\frac1{2\varepsilon}\mathbb E\bigl[\mathbf 1_{|f(t)|<\varepsilon}|f'(t)|\bigr]&=\frac1{2\varepsilon}\int_{-\varepsilon}^\varepsilon\frac{e^{-y^2/2A}}{\sqrt{2\pi A}}\,\mathbb E\Bigl|\frac{By}A+\sqrt{C-\tfrac{B^2}A}\,Z\Bigr|\dd y\\
&\longrightarrow\frac1{\sqrt{2\pi A}}\sqrt{\frac2\pi}\sqrt{C-\frac{B^2}A}=\frac{\sqrt{AC-B^2}}{\pi A},
\end{aligned}
\]
and the left side is bounded by $\frac{|B|+\sqrt C}{\sqrt{2\pi}}$ for $\varepsilon\le1$ (as $A\ge1$), which is bounded on $[a,b]$; dominated convergence in $t$ gives the formula.

(d) $A'=2B$ and $(tA')'=\sum_j4j^2t^{2j-1}=4tC$. Hence
$\frac1{4t}\bigl(\frac{tA'}A\bigr)'=\frac1{4t}\bigl(\frac{(tA')'}A-\frac{tA'^2}{A^2}\bigr)=\frac CA-\frac{B^2}{A^2}=\pi^2\rho_n^2$. With $A=\frac{1-t^{2n+2}}{1-t^2}$,
$\frac{tA'}A=\frac{2t^2}{1-t^2}-\frac{(2n+2)t^{2n+2}}{1-t^{2n+2}}$; differentiating and dividing by $4t$ gives $\frac1{(1-t^2)^2}-\frac{(n+1)^2t^{2n}}{(1-t^{2n+2})^2}$.

(e) $t^nf(\frac1t)=\sum_ja_{n-j}t^j$ has the same law as $f$ and its roots are the inverses of those of $f$; and $f(-t)=\sum_j(-1)^ja_jt^j$ has the same law as $f$. So the
four counts have the same law, and since $0$, $1$ and $-1$ are almost surely not roots, $\mathbb EN_n=4\int_0^1\rho_n$.

\emph{Upper bound.} $\rho_n(t)\le\frac1{\pi(1-t^2)}$, and $\int_0^{1-1/n}\frac{\dd t}{\pi(1-t^2)}=\frac1{2\pi}\log(2n-1)$. On $[1-\frac1n,1]$, $\rho_n\le\frac1\pi\sqrt{C/A}$ with
$C\le n^3$ and $A\ge(n+1)t^{2n}\ge\frac{n+1}{16}$ (as $(1-\frac1n)^n\ge\frac14$ for $n\ge2$), so $\rho_n\le\frac{4n}\pi$ there and the integral is at most $\frac4\pi$. So
$\mathbb EN_n\le\frac2\pi\log(2n)+\frac{16}\pi$.

\emph{Lower bound.} Let $\delta=\frac{\log n}n$ and $0\le t\le1-\delta$, $s=1-t$. Then $t^{2n}\le e^{-2ns}$ and $(1-t^2)^2\le4s^2$, and $s^2e^{-2ns}$ decreases for $s\ge\frac1n$, so
$(n+1)^2t^{2n}(1-t^2)^2\le4(n+1)^2\delta^2n^{-2}\le\frac{16\log^2n}{n^2}=\eta_n$, while $(1-t^{2n+2})^2\ge\frac12$. So the second term in Kac's density is at most $2\eta_n$ times the first,
$\rho_n(t)\ge\frac{\sqrt{1-2\eta_n}}{\pi(1-t^2)}$, and $\int_0^{1-\delta}\rho_n\ge\frac{\sqrt{1-2\eta_n}}{2\pi}\log\frac1\delta$. Hence $\mathbb EN_n\ge\frac2\pi\sqrt{1-2\eta_n}(\log n-\log\log n)$, and
$\mathbb EN_n\sim\frac2\pi\log n$.

\begin{figure}[ht]
\centering\fig{ch01-kac}
\caption{Kac's density $\rho_n$ on $[0,1]$ for $n=10$ and $n=100$, with the envelope $\frac1{\pi(1-t^2)}$ (dashed). The density follows the envelope until $1-t$ is of order
$\frac1n$, and then climbs to $\rho_n(1)=\frac1\pi\sqrt{\frac{n(n+2)}{12}}$, which is $9.3$ for $n=100$ (off the picture, arrow): the real roots crowd near $\pm1$.}
\label{fig:kac}
\end{figure}

(f) Numerically, one finds $\mathbb EN_{100}=3.56$. The other roots are complex, and numerical evaluation shows where: almost all of them lie close to the unit circle, $81\%$ within
$0.1$ of it for $n=50$, spread almost uniformly in angle. Real roots are rare because the real axis meets that crowded annulus only near $\pm1$.

\tools{Volume~4, Chapter~6 (Gaussian vectors); Volume~7, Chapters~8 and~9 (dominated convergence, Fubini).}
\leads Kac found the formula in 1943 and the asymptotics $\frac2\pi\log n$; numerical evaluation shows $\mathbb EN_n-\frac2\pi\log n\to0.6257\ldots$. The same counting formula is the
\emph{Kac--Rice formula}, a basic tool for the zeros of random functions and the level crossings of Gaussian processes. For random polynomials, Ibragimov and Maslova (1971)
showed that the answer $\frac2\pi\log n$ holds for all centred coefficients in the domain of attraction of the normal law, and Tao and Vu (2015) studied the universality of
local statistics of the roots.
\end{solution}

\section*{Map of the chapter}

\begin{center}
\small
\adjustbox{max width=\linewidth}{\begin{tabular}{lll}
Problem & Topic & Volumes used\\\hline
\ref{pb:pairwisebits} & pairwise independent signs: Chebyshev yes, Gauss no & 3\\
\ref{pb:sawtooth} & Kac's lacunary sums $\sum B(2^kx)$ & 4, 10\\
\ref{pb:peaks} & peaks of permutations, $m$-dependent central limit theorem & 3\\
\ref{pb:tilings} & dominoes in a random tiling, renewal--reward & 3, 10\\
\ref{pb:harper} & Harper's theorem on the blocks of a random partition & 0, 3, 4\\
\ref{pb:fixedtv} & fixed points: exact distance to Poisson & 3, 10\\
\ref{pb:kubslow} & Stein's bound for the Kubilius model, and a lower bound & 5, 10\\
\ref{pb:twelve} & a fourth-order Lindeberg swap; twelve uniforms & 3\\
\ref{pb:threeap} & three-term progressions in a random set (Chen--Stein) & 10\\
\ref{pb:spearman} & Spearman's $\rho$ by exchangeable pairs & 3, 10\\
\ref{pb:erdostaylor} & the Erdős--Taylor theorem on planar returns & 3, 4, 6\\
\ref{pb:dice} & dice modulo $q$; a local limit theorem with span $3$ & 2, 10\\
\ref{pb:squaresmodp} & sums of squares modulo $p$ mix in two steps & 1, 9\\
\ref{pb:lucky} & lucky tickets and the Edgeworth correction & 3, 4, 10\\
\ref{pb:signedsums} & $\pm1\pm2\cdots\pm n=0$ & 4, 10\\
\ref{pb:smirnov} & Smirnov's two-sample theorem; theta identity & 3, 4, 10\\
\ref{pb:primeseries} & random series over the primes & 3, 5, 10\\
\ref{pb:factorialseries} & $\sum\varepsilon_n/n!$, a law of dimension $0$ & 6\\
\ref{pb:smooth} & $\sum\varepsilon_n/n^2$ has a smooth density & 7, 10\\
\ref{pb:purity} & the Jessen--Wintner purity theorem & 6, 10\\
\ref{pb:golden} & Erdős's theorem on the golden Bernoulli convolution & 7, 12\\
\ref{pb:kaczeros} & Kac's formula for the real roots of a random polynomial & 3, 6\\
\end{tabular}}
\end{center}

\section*{Sources and further reading}

\begin{itemize}
\item H.~Cramér, Half a century with probability theory: some personal recollections, \emph{Annals of Probability} 4 (1976), 509--546 (the opening story);
J.~W.~Lindeberg, Eine neue Herleitung des Exponentialgesetzes in der Wahrscheinlichkeitsrechnung, \emph{Mathematische Zeitschrift} 15 (1922), 211--225.
\item M.~Kac, On the distribution of values of sums of the type $\sum f(2^kt)$, \emph{Annals of Mathematics} 47 (1946), 33--49.
\item L.~H.~Harper, Stirling behavior is asymptotically normal, \emph{Annals of Mathematical Statistics} 38 (1967), 410--414; J.~N.~Darroch, On the distribution of the number of
successes in independent trials, \emph{Annals of Mathematical Statistics} 35 (1964), 1317--1321; J.~Pitman, Probabilistic bounds on the coefficients of polynomials with only real
zeros, \emph{Journal of Combinatorial Theory, Series A} 77 (1997), 279--303.
\item A.~D.~Barbour, L.~Holst and S.~Janson, \emph{Poisson Approximation}, Oxford University Press, 1992; L.~H.~Y.~Chen, L.~Goldstein and Q.-M.~Shao, \emph{Normal Approximation by
Stein's Method}, Springer, 2011; C.~Stein, \emph{Approximate Computation of Expectations}, IMS Lecture Notes 7, 1986.
\item W.~Hoeffding, A combinatorial central limit theorem, \emph{Annals of Mathematical Statistics} 22 (1951), 558--566; E.~Bolthausen, An estimate of the remainder in a combinatorial
central limit theorem, \emph{Zeitschrift für Wahrscheinlichkeitstheorie und verwandte Gebiete} 66 (1984), 379--386.
\item A.~Rényi and P.~Turán, On a theorem of Erdős--Kac, \emph{Acta Arithmetica} 4 (1958), 71--84.
\item G.~E.~P.~Box and M.~E.~Muller, A note on the generation of random normal deviates, \emph{Annals of Mathematical Statistics} 29 (1958), 610--611.
\item D.~Conlon and W.~T.~Gowers, Combinatorial theorems in sparse random sets, \emph{Annals of Mathematics} 184 (2016), 367--454; M.~Schacht, Extremal results for random discrete
structures, \emph{Annals of Mathematics} 184 (2016), 333--365.
\item P.~Erdős and S.~J.~Taylor, Some problems concerning the structure of random walk paths, \emph{Acta Mathematica Academiae Scientiarum Hungaricae} 11 (1960), 137--162.
\item P.~Diaconis, \emph{Group Representations in Probability and Statistics}, IMS Lecture Notes 11, 1988 (the upper bound lemma and random walks modulo $p$).
\item J.~H.~van Lint, Representation of $0$ as $\sum_{k=-N}^N\varepsilon_kk$, \emph{Proceedings of the American Mathematical Society} 18 (1967), 182--184; R.~C.~Entringer, Representation
of $m$ as $\sum_{k=-n}^n\varepsilon_kk$, \emph{Canadian Mathematical Bulletin} 11 (1968), 289--293.
\item A.~N.~Kolmogorov, Sulla determinazione empirica di una legge di distribuzione, \emph{Giornale dell'Istituto Italiano degli Attuari} 4 (1933), 83--91; N.~V.~Smirnov, On the
estimation of the discrepancy between empirical curves of distribution for two independent samples, \emph{Bulletin Mathématique de l'Université de Moscou} 2 (1939), 3--14;
B.~V.~Gnedenko and V.~S.~Korolyuk, On the maximum discrepancy between two empirical distributions, \emph{Doklady Akademii Nauk SSSR} 80 (1951), 525--528.
\item H.~Bohr and B.~Jessen, Über die Werteverteilung der Riemannschen Zetafunktion, \emph{Acta Mathematica} 54 (1930), 1--35; B.~Jessen and A.~Wintner, Distribution functions
and the Riemann zeta function, \emph{Transactions of the American Mathematical Society} 38 (1935), 48--88.
\item P.~Erdős, On a family of symmetric Bernoulli convolutions, \emph{American Journal of Mathematics} 61 (1939), 974--976; R.~Salem, A remarkable class of algebraic integers,
\emph{Duke Mathematical Journal} 11 (1944), 103--108; B.~Solomyak, On the random series $\sum\pm\lambda^n$ (an Erdős problem), \emph{Annals of Mathematics} 142 (1995), 611--625;
P.~Shmerkin, On the exceptional set for absolute continuity of Bernoulli convolutions, \emph{Geometric and Functional Analysis} 24 (2014), 946--958; P.~P.~Varjú, Absolute continuity
of Bernoulli convolutions for algebraic parameters, \emph{Journal of the American Mathematical Society} 32 (2019), 351--397.
\item M.~Kac, On the average number of real roots of a random algebraic equation, \emph{Bulletin of the American Mathematical Society} 49 (1943), 314--320; A.~Edelman and E.~Kostlan,
How many zeros of a random polynomial are real?, \emph{Bulletin of the American Mathematical Society} 32 (1995), 1--37; I.~A.~Ibragimov and N.~B.~Maslova, The mean number of real
zeros of random polynomials I, \emph{Theory of Probability and its Applications} 16 (1971), 228--248; T.~Tao and V.~Vu, Local universality of zeroes of random polynomials,
\emph{International Mathematics Research Notices} 2015, 5053--5139.
\end{itemize}
